\subsection{可解圈空间}

定义 2. --- 称拓扑空间 B 为可解圈的，如果它是局部道路连通的，并且 B 的每一点 b 都有一个邻域 V，使得由典范单射导出的、从 \(\pi_1(V,b)\) 到 \(\pi_1(B,b)\) 中的同态以单位元为其像。

一个局部道路连通的拓扑空间是可解圈的，当且仅当它的每个连通分支都是可解圈的。

设 B 是一个局部道路连通的拓扑空间。假设 B 的每一点都有一个可解圈的邻域 V。那么 B 是可解圈的。

每个局部道路连通且道路单连通的拓扑空间都是可解圈的。同伦于一点的局部道路连通空间尤其如此。

注. --- 1) 存在连通且局部道路连通的拓扑空间 B，使得 B 的某些点 a（但并非所有点）具有邻域 V，使从 \(\pi_{1}(V,a)\) 到 \(\pi_{1}(B,a)\) 中的同态是平凡的（III, p. 336, 习题 6）。

\begin{enumerate}
\def\labelenumi{\arabic{enumi})}
\setcounter{enumi}{1}
\tightlist
\item
  设 B 是一个可解圈空间；群胚 \(\varpi(\mathrm{B})\) 在 B-空间上的每个作用都无局部单值（参见 III, p. 313, 注）。特别地，为使映射 \(p: E \rightarrow B\) 使 E 成为 B 的覆叠，必要且充分的是：它平展、分离，并且满足道路提升性质（III, p. 315, 推论 3）。
\end{enumerate}

命题 1. --- 有限多个可解圈空间的乘积空间是可解圈的。

只需证明：若 A 与 B 是两个可解圈空间，则其乘积 A × B 也是。在这些条件下，空间 A × B 确实是局部道路连通的（III, p. 261, 命题 9）。设 (a, b) 是 A × B 的一点。依假设，存在 a 的一个邻域 U（相应地，b 的一个邻域 V），使同态 \(i_* \colon \pi_1(\mathrm{U}, a) \to \pi_1(\mathrm{A}, a)\)（它由典范单射 \(i \colon \mathrm{U} \to \mathrm{A}\) 导出；相应地，同态 \(j_* \colon \pi_1(\mathrm{V}, b) \to \pi_1(\mathrm{B}, b)\) 由典范单射 \(j \colon \mathrm{V} \to \mathrm{B}\) 导出）的像退缩为单位元。于是 U × V 是 (a, b) 在 A × B 中的一个邻域。同态 \(\pi_1(\mathrm{U} \times \mathrm{V}, (a, b)) \to \pi_1(\mathrm{A} \times \mathrm{B}, (a, b))\) 与同态 \((i_*, j_*)\) 等同（III, p. 297, 命题 4 的推论）。因此其像退缩为单位元。这就证明了 A × B 是可解圈的。

命题 2. --- a) 可解圈空间的每个覆叠都是可解圈的。

\begin{enumerate}
\def\labelenumi{\alph{enumi})}
\setcounter{enumi}{1}
\item
  反之，设 \(p \colon \mathrm{E} \to \mathrm{B}\) 是一个平展且满射的映射，并设 \(\mathrm{E}\) 是可解圈空间。那么 \(\mathrm{B}\) 是可解圈的。
\item
  设 B 是一个可解圈拓扑空间，\((\mathrm{E}, p)\) 是 B 的一个覆叠。于是空间 E 局部道路连通。设 x 是 E 的一点，记 \(b = p(x)\)，并设 V 是 b 在 B 中的一个邻域，使典范同态 \(\pi_{1}(\mathrm{V}, b) \to \pi_{1}(\mathrm{B}, b)\) 平凡。用 \(i: V \to B\) 与 \(j: p^{-1}(V) \to E\) 表示典范单射。依假设，同态 \(\pi_{1}(i, b)\) 的像退缩为单位元。由于 \(p \circ i = j \circ p\) 且 \(\pi_{1}(p, x)\) 是单射的，同态 \(\pi_{1}(j, x)\) 的像退缩为单位元。这就证明了 E 是可解圈的。
\item
  空间 B 局部道路连通。设 b 是 B 的一点，x 是 \(E_{b}\) 的一点。由于 p 平展，存在 x 在 E 中的一个邻域 U，使 p 诱导 U 到 \(p(\mathrm{U})\) 上的一个同胚。设 V 是 x 在 E 中的一个包含于 U 的邻域，使同态 \(\pi_{1}(j,x)\) 平凡，其中 \(j:V\to E\) 表示典范单射。用 \(q:V\to p(V)\) 表示由 p 取子空间而导出的映射；它是一个同胚。再用 i 表示 \(p(\mathrm{V})\) 到 B 中的典范单射。我们有 \(p\circ j=i\circ q\)，因此 \(\pi_{1}(i,b)\circ\pi_{1}(q,x)=\pi_{1}(p,x)\circ\pi_{1}(j,x)\) 是平凡同态。由于同态 \(\pi_{1}(q,x)\) 是同构，同态 \(\pi_{1}(i,b)\) 平凡。由此得出空间 B 是可解圈的。
\end{enumerate}

命题 3. --- 设 B 是一个可解圈拓扑空间。设 (Y, q) 是 B 的一个覆叠，(Z, p) 是 Y 的一个覆叠。于是赋以映射 \(q \circ p\) 的拓扑空间 Z 是 B 的一个覆叠。

事实上，映射 \(q \circ p\) 是平展的（I, p. 29, 命题 6）且分离的（I, p. 25, 命题 2）。由上面的注 2，因此只需证明它满足道路提升性质。设 \(z\) 是 Z 的一点，并设 \(c \colon \mathbf{I} \to \mathbf{B}\) 是 B 中以 \(q(p(z))\) 为起点的一条道路。存在一条道路 \(c'\)，它以 \(p(z)\) 为起点、位于 Y 中，且提升 \(c\)，因为 Y 是 B 的覆叠（III, p. 302, 推论 2）。由于 Z 是 Y 的覆叠，存在一条道路 \(c''\)，它以 \(z\) 为起点、位于 Z 中，且提升 \(c'\)；道路 \(c''\) 提升 \(c\)。命题得证。

\subsection{可解圈空间的万有覆叠}

设 B 是一个拓扑空间，b 是 B 的一点。回忆 \(\Lambda_{b}(\mathrm{B})\) 表示 \(\mathcal{C}_{c}(\mathbf{I};\mathrm{B})\) 中由以 b 为起点的道路组成的子空间，赋以紧收敛拓扑的商拓扑。我们用 \(e_{\mathrm{B}}\colon\Lambda_{b}(\mathrm{B})\to\mathrm{B}\) 表示把一条道路对应到其终点的映射；它是连续的。用 \(\lambda_{b}(\mathrm{B})\) 表示 \(\Lambda_{b}(\mathrm{B})\) 关于严格同伦关系的商空间，并赋以商拓扑。由于两条严格同伦的道路有相同的终点，映射 \(e_{B}\) 通过过渡到商，定义一个连续映射 \(\varepsilon_{\mathrm{B}}\colon\lambda_{b}(\mathrm{B})\to\mathrm{B}\)，称为终点映射。

定理 1. --- 设 B 是一个连通且局部道路连通的拓扑空间，并设 b 是 B 的一点。下列性质等价：

\begin{enumerate}
\def\labelenumi{(\roman{enumi})}
\item
  空间 B 是可解圈的。
\item
  存在 B 的一个非空且道路单连通的覆叠。
\item
  空间 \(\lambda_{b}(\mathrm{B})\)，赋以终点映射 \(\varepsilon_{\mathrm{B}}\colon\lambda_{b}(\mathrm{B})\to\mathrm{B}\)，是 B 的一个覆叠。
\item
  群 \(\pi_1(B, b)\) 对容许拓扑是离散的。
\item
  群 \(\pi_{1}(B,b)\) 对紧收敛拓扑的商拓扑是离散的。
\end{enumerate}

此外，当这些条件满足时，空间 \(\lambda_{b}(\mathrm{B})\) 是道路单连通的、单连通的、以 \(\pi_{1}(\mathrm{B},b)^{\circ}\) 为群的伽罗瓦覆叠，且点覆叠 \((\lambda_{b}(\mathrm{B}),\varepsilon_{b})\) 是点空间 \((\mathrm{B},b)\) 的一个万有覆叠。

由可解圈空间的定义可知，\(\pi_{1}(B,b)\) 中退缩为单位元的子群是容许的，当且仅当 B 是可解圈的。这就证明了 (i)⇔(iv)。此外，性质 (iv) 与 (v) 的等价性由 III, p. 320 的注 4 与注 5 得出。

(iv)⇒(ii)。由 III, p. 316, 命题 5，存在 B 的一个覆叠 (E, p) 及一点 \(x \in E_{b}\)，使 \(p_{*}(\pi_{1}(E, x)) = \{e\}\)；特别地，\(\pi_{1}(E, x)\) 退缩为单位元。于是 x 在 E 中的弧连通分支是 B 的一个非空覆叠（I, p. 80, 命题 6 的推论 1），且是道路单连通的。

(ii)⇒(iii)。设 E 是 B 的一个非空道路单连通覆叠。设 x 是 \(E_{b}\) 的一点（这样的点存在，因为 B 连通，I, p. 74, 命题 4）。投影 \(q: E \to B\) 诱导一个同胚 \(\Lambda_{x}(E) \to \Lambda_{b}(B)\)（III, p. 302, 命题 3 的推论 2），从而通过过渡到弧连通分支，得到一个同胚 \(q_{*}: \lambda_{x}(E) \simeq \lambda_{b}(B)\)。终点映射 \(\Lambda_{x}(E) \to E\) 是连续且开的，因为 E 局部道路连通（III, p. 264, 推论）。因此它通过过渡到弧连通分支所定义的映射 \(\varepsilon_{\mathrm{E}}: \lambda_{x}(E) \to E\) 是连续且开的。映射 \(\varepsilon_{\mathrm{E}} \circ (q_{*})^{-1}: \lambda_{b}(B) \to E\) 是双射的、连续的且开的。这就证明了赋以紧收敛拓扑与终点映射的拓扑空间 \(\lambda_{b}(B)\) 是 B 的一个覆叠。

(iii)⇒(v)。事实上，赋以紧收敛拓扑的商拓扑的集合 \(\pi_{1}(B,b)\)，与 B-空间 \(\lambda_{b}(B)\) 在 b 上、在 b 处常值圈所属的那个类处的纤维等同。若 \(\lambda_{b}(B)\) 是 B 的覆叠，则 \(\pi_{1}(B,b)\) 对紧收敛拓扑是离散的。

假设这些断言成立。由前文，空间 \(\lambda_b(B)\) 是 B 的一个道路单连通的覆叠，因而是单连通的。于是点空间 ( \(\lambda_b(B), \varepsilon_b\) ) 是 (B, \(b\) ) 的一个万有覆叠（I, p. 126, 命题 3 的推论），且覆叠 \(\lambda_b(B)\) 是 B 的一个伽罗瓦覆叠（I, p. 120, 命题 1）。典范同态 \(h_{(\lambda_b(B), \varepsilon_b)}: \pi_1(B, b) \to \text{Aut}_B(\lambda_b(B))\) 于是是一个同构（III, p. 306, 命题 5）。

注 1. --- 设 B 是一个连通可解圈空间，并设 b 是 B 的一点。由定理 1, (iv) 可知，群 \(\pi_{1}(B,b)\) 的每个作用都是容许的。因此，每个 \(\pi_{1}(B,b)\) -集合都同构于一个 \(\pi_{1}(B,b)\) -集合 \(E_{b}\)，其中 E 是 B 的覆叠（III, p. 318, 命题 7）。还回忆（III, p. 310, 命题 2）：若 E 与 \(E'\) 是 B 的覆叠，则映射 \(f \mapsto f_{b}\) 诱导一个从 E 到 \(E'\) 中的 B-空间态射所成之集到 \(\pi_{1}(B,b)\) -集合的从 \(E_{b}\) 到 \(E'_{b}\) 上的态射所成之集上的双射。

用范畴的语言来说：把 B 的覆叠 E 对应到纤维 \(E_b\) 的函子，是从 B 的覆叠的范畴到 \(\pi_1(B,b)\) -集合的范畴的一个范畴等价。*

推论 1. --- 设 B 是一个可解圈拓扑空间，并设 b 是 B 的一点。设 E 是 B 的一个连通覆叠，x 是纤维 \(E_{b}\) 的一点。下列性质等价：

\begin{enumerate}
\def\labelenumi{(\roman{enumi})}
\item
  群 \(\pi_{1}(\mathrm{E},x)\) 是平凡的。
\item
  空间 E 是单连通的。
\item
  点空间 (E, x) 是 (B, b) 的一个万有覆叠。
\end{enumerate}

蕴涵关系 (i)⇒(ii) 已仅在“空间 B 连通且局部道路连通”这一假设下证明（III, p. 309, 命题 1 的推论 2），而蕴涵关系 (ii)⇒(iii) 则在没有任何假设下证明（I, p. 126, 推论）。

我们来证明 (iii)⇒(i)。由 IV, p. 342 的定理 1，存在 B 的一个覆叠 E' 及一点 \(x'\)（属于纤维 \(E'_{b}\)），使群 \(\pi_1(E', x')\) 退缩为单位元。在假设 (iii) 下，存在

一个点态 B-态射 \(f \colon (\mathrm{E}, x) \to (\mathrm{E}', x')\)。用 \(p \colon \mathrm{E} \to \mathrm{B}\) 与 \(p' \colon \mathrm{E}' \to \mathrm{B}\) 表示覆叠 \(\mathrm{E}\) 与 \(\mathrm{E}'\) 的投影；我们有 \(p = p' \circ f\)，因此 \(\pi_1(p, x) = \pi_1(p', x') \circ \pi_1(f, x)\)。由于群 \(\pi_1(\mathrm{E}', x')\) 平凡，单射同态 \(\pi_1(p, x)\) 以单位元为其像。这就证明了群 \(\pi_1(\mathrm{E}, x)\) 退缩为单位元。

推论 2. --- 设 B 是一个可解圈拓扑空间，并设 b 是 B 的一点。设 (E, x) 是点空间 (B, b) 的一个万有覆叠。对 B 的每个覆叠 E' 及每一点 \(x' \in E'_{b}\)，使 f(x) = x' 的唯一的 B-态射 f: E → E' 使 E 成为 E' 的一个覆叠。于是点空间 (E, x) 是 (E', x') 的一个万有覆叠。

由 I, p. 81 的命题 7，赋以映射 f 的空间 E 是 E' 的一个覆叠。最后的断言于是由推论 1 得出。

设 B 是一个可解圈拓扑空间，并设 \(a\)、\(b\) 是 B 的两点。商拓扑空间 \(\varpi_{a,b}(\mathrm{B})\) 同胚于赋以紧收敛拓扑的商拓扑的空间 \(\pi_1(\mathrm{B},a)\)（III, p. 293, 注 3），因而由 IV, p. 342 的定理 1 是离散的。所以，B 中连接 \(a\) 与 \(b\) 的每条道路的严格同伦类都是 \(\Lambda_{a,b}(\mathrm{B})\) 的一个开子集。更一般地：

命题 4. --- 设 B 是一个可解圈拓扑空间。严格同伦是拓扑空间 \(\mathcal{C}_{c}(\mathbf{I};\mathbf{B})\) 中的一个开等价关系。

首先假设空间 B 是道路单连通的。设 \(\varphi\colon\Lambda(B)\to B\times B\) 是把 B 中的道路 \(c\) 对应到由其起点与终点组成的偶 \((c(0),c(1))\) 的映射。为使 B 中的两条道路严格同伦，必要且充分的是它们有相同的起点与相同的终点（III, p. 292, 命题 2）。由于空间 B 连通且局部道路连通，映射 \(\varphi\) 是满射的、连续的且开的（III, p. 262, 命题 10）。因此严格同伦关系是开的（TG, I, p. 32, 命题 3）。

在一般情形，设 E 是 B 的一个道路单连通覆叠，当 \(B \neq \varnothing\) 时非空。由于 B 连通，这个覆叠是满射的（I, p. 74, 命题 4）。用 p 表示 B-空间 E 的投影，并用 \(\widetilde{p} \colon \Lambda(\mathrm{E}) \to \Lambda(\mathrm{B})\) 表示把 E 中的道路 c 对应到道路 \(p \circ c\) 的映射。赋以映射 \(\widetilde{p}\) 的 \(\Lambda(\mathrm{E})\) 是 \(\Lambda(\mathrm{B})\) 的一个覆叠（III, p. 302, 命题 3 的推论 1）。为使 B 中的两条道路严格同伦，必要且充分的是它们都是 E 中两条严格同伦的道路在映射 \(\widetilde{p}\) 下的像（III, p. 302, 命题 4）。设 U 是空间 \(\Lambda(B)\) 的一个开子集。集合 \((\widetilde{p})(U)\) 在 \(\Lambda(E)\) 中开；由证明的第一部分，集合 \(U'\)（关于严格同伦关系、相对于 \((\widetilde{p})(U)\) 饱和）在 \(\Lambda(E)\) 中开。由于 \(\Lambda(E)\) 是 \(\Lambda(B)\) 的覆叠，集合 \(\widetilde{p}(U')\) 在 \(\Lambda(B)\) 中开。由于 \(\widetilde{p}(U')\) 是 U 关于严格同伦关系的饱和化，命题得证。

\subsection{例子}

\subitemhead{局部可缩空间}

定义 3. --- 称拓扑空间 B 为局部可缩的，如果 B 的每一点 b 都有一个邻域 V，使点空间 (V, b) 是可缩的（III, p. 234, 例 3）。

命题 5. --- 局部可缩的拓扑空间是可解圈的。

设 B 是一个局部可缩空间。设 b 是 B 的一点，V 是 b 的一个邻域，使 \((V, b)\) 是一个点态可缩空间。空间 V 同伦等价于一点，故 \(\pi_{1}(V, b) = \{\varepsilon_{b}\}\)（III, p. 296, 推论 3）。

为证明本命题，还需证明空间 B 局部道路连通。设 \(\sigma: V \times I \to V\) 是 \(b\) 处的一个连接像为 \(b\) 的常值映射与映射 \(Id_V\) 的点态同伦。对 \(b\) 的每个包含于 V 的邻域 W，集合 \(\sigma(W \times I)\) 是 \(b\) 的一个邻域，因为它包含 \(W = \sigma(W \times \{1\})\)；而且它道路连通，因为对每个 \((a, s) \in W \times I\)，映射 \(t \mapsto \sigma(a, ts)\) 是连接 \(b = \sigma(a, 0)\) 与 \(\sigma(a, s)\) 的一条 \(\sigma(W \times I)\) 中的道路。我们来证明诸形如 \(\sigma(W \times I)\) 的集合构成 \(b\) 的一个基本邻域系。设 \(V'\) 是 \(b\) 在 B 中的一个开邻域；于是 \(\sigma^{-1}(V')\) 是 \(\{b\} \times I\) 在 \(V \times I\) 中的一个开邻域。由于空间 I 紧，投影 \(\mathrm{pr}_1: V \times I \to V\) 是逆紧的（TG, I, p. 77, 推论 5），且 \(\mathrm{pr}_1(\complement\,\sigma^{-1}(V'))\) 是 V 的一个不包含 \(b\) 的闭子集。于是它的补 W 是 \(b\) 在 V 中的一个开邻域，使 \(\mathrm{W} \times \mathbf{I} \subset \sigma^{-1}(\mathrm{V}')\)，从而 \(\sigma(\mathrm{W} \times \mathbf{I}) \subset \mathrm{V}'\)。

数值空间或实、复 n 维射影空间的每个开子集（参见第六章与第八章）都是可解圈的，欧氏球面 \(S_{n}\)（TG, VI, p. 11, 命题 4）亦然。更一般地，每个拓扑流形（VAR R, 第二部分, p. 7, 记号与约定）都是可解圈的。事实上，这些空间都是局部可缩的。

\subitemhead{圆的庞加莱群}

拓扑空间 R 同伦于一点（III, p. 234, 例 4），因此是道路单连通的（IV, p. 340）。R 到 T = R/Z 上的典范满射 p 使之成为以 Z 为群的主覆叠（I, p. 100, 例 4）。拓扑空间 T 是可解圈的，且 (R, 0) 是 (T, 0) 的一个万有覆叠。由此我们导出群的一个典范同构 \(\pi_{1}(\mathbf{T}, 0) \to \mathbf{Z}\)：它把 0 处的圈 \(t \mapsto p(nt)\) 的类对应到 Z 的元素 n。顾及例 6（I, p. 101），我们导出下列命题。

命题 6. --- 映射 \(p \colon x \mapsto e^{2\pi ix}\)（从 \(\mathbf{R}\) 到 \(\mathbf{S}_1\) 上）使 \((\mathbf{R}, 0)\) 成为 \((\mathbf{S}_1, 1)\) 的一个万有覆叠。群 \(\pi_1(\mathbf{S}_1, 1)\) 同构于 \(\mathbf{Z}\)；圈 \(t \mapsto e^{2\pi it}\) 的类是一个生成元。对每个整数 \(n > 0\)，映射 \(z \mapsto z^n\) 使 \(\mathbf{S}_1\) 成为一个以 \(\mathbf{Z}/n\mathbf{Z}\) 为群的、\(\mathbf{S}_1\) 上的伽罗瓦覆叠。\(\mathbf{S}_1\) 的每个连通且非空的覆叠都同构于这些覆叠之一。

只需证明最后的断言。一个连通且非空的覆叠 E 在 1 上的纤维 \(E_{1}\)（E 为 \(\mathbf{S}_{1}\) 的覆叠）赋有群 \(\pi_{1}(\mathbf{S}_{1},1)\) 的一个传递作用。由于 Z 的子群都是形如 nZ 的集合（对 \(n \geqslant 0\)），可见 \(E_{1}\) 同构于与前述诸覆叠相伴的齐性集合之一。由于 \(\mathbf{S}_{1}\) 连通且局部道路连通，E 同构于这些覆叠之一。

推论. --- 映射 \(z \mapsto e^{z}\)（从 C 到 \(C^{*}\)）使 \((\mathbf{C},0)\) 成为 \((\mathbf{C}^{*},1)\) 的一个万有覆叠。群 \(\pi_{1}(\mathbf{C}^{*},1)\) 同构于 Z；圈 \(t \mapsto e^{2\pi it}\) 的类是一个生成元。对每个整数 n \textgreater{} 0，映射 \(z \mapsto z^{n}\) 使 \(C^{*}\) 成为一个以 Z/nZ 为群的、\(C^{*}\) 上的伽罗瓦覆叠。\(C^{*}\) 的每个连通且非空的覆叠都同构于这些覆叠之一。

映射 \(z \mapsto (|z|, z/|z|)\) 是空间 \(\mathbf{C}^*\) 到空间 \(\mathbf{R}_+^* \times \mathbf{S}^1\) 上的一个同胚（TG, VI, p. 10, 命题 3）；特别地，\(\mathbf{C}^*\) 是道路连通的。由于映射 \(x \mapsto e^x\)（从 \(\mathbf{R}\) 到 \(\mathbf{R}_+^*\) 上）是一个同胚，由命题 6 可知，映射 \((x, y) \mapsto e^x e^{2\pi iy}\) 使 \(\mathbf{R}^2\) 成为 \(\mathbf{C}^*\) 的一个覆叠。把 \(\mathbf{R}^2\) 与 \(\mathbf{C}\) 依映射 \((x, y) \mapsto x + iy\) 等同起来，我们断定：空间 \(\mathbf{C}\)——赋以映射 \(z \mapsto e^z\)——是 \(\mathbf{C}^*\) 的一个覆叠。由于空间 \(\mathbf{C}\) 连通且道路单连通，点覆叠 \((\mathbf{C}, 0)\) 是点空间 \((\mathbf{C}^*, 1)\) 的一个万有覆叠。

由该推论（III, p. 297）还可得出：\(\mathbf{C}^*\) 在 1 处的庞加莱群同构于 \(\mathbf{Z}\)，且类 \(\gamma\)（圈 \(t \mapsto e^{2\pi it}\) 的类）是一个生成元。

  设 n 是整数 \textgreater{} 0。映射 \(x \mapsto x^{n}\)（从 \(R_{+}^{*}\) 到其自身）是一个同胚；我们同前推出：映射 \(z \mapsto z^{n}\) 使 \(C^{*}\) 成为 \(C^{*}\) 的一个 n 次覆叠。它是以 Z/nZ 为群的伽罗瓦覆叠。最后的断言如命题 6 的证明中那样证明。

对每个整数 \(n \geqslant 0\)，环面 \((\mathbf{S}_{1})^{n}\) 是可解圈空间（IV, p. 341, 命题 1），且它在每一点的庞加莱群都同构于 \(\mathbf{Z}^{n}\)（III, p. 297, 命题 4）。

我们将在第 3 节中推广这些结果，该节专门讨论拓扑群的覆叠。

\subitemhead{实射影空间}

设 n 是整数 \textgreater{} 0。空间 \(S_{n}\) 与 \(\mathbf{P}_{n}(\mathbf{R})\) 都是连通的、非空的，且典范映射（TG, VI, p. 13）\(\varphi\colon\mathbf{S}_{n}\to\mathbf{P}_{n}(\mathbf{R})\) 使 \(S_{n}\) 成为 \(\mathbf{P}_{n}(\mathbf{R})\) 的一个主覆叠，其群为 \(\{+1,-1\}\)（I, p. 99, 例 1）。对 \(n\geqslant2\)，球面 \(S_{n}\) 是单连通的（I, p. 127, 例 3）且是可解圈的，因而是道路单连通的（IV, p. 344, 推论 1）。对 n=1，映射 \(z\mapsto z^{2}\)（从 \(S_{1}\) 到 \(S_{1}\)）在 \(S_{1}\) 中定义了一个等价关系，其诸等价类是 \(S_{1}\) 的对径点偶。由过渡到商所定义的映射 \(\varphi\colon\mathbf{S}_{1}\to\mathbf{P}_{1}(\mathbf{R})\) 定义了一个连续双射 \(\psi\colon\mathbf{S}_{1}\to\mathbf{P}_{1}(\mathbf{R})\)，使 \(\psi(z^{2})=\varphi(z)\)。由于 \(S_{1}\) 是紧空间，映射 \(\psi\) 是一个同胚（TG, I, p. 63, 推论 2）。

命题 7. --- 映射 \(x \mapsto \varphi(e^{2\pi ix})\)（从 \(\mathbf{R}\) 到 \(\mathbf{P}_{1}(\mathbf{R})\) 上）使 \((\mathbf{R},0)\) 成为 \((\mathbf{P}_{1}(\mathbf{R}),\varphi(1))\) 的一个万有覆叠。设 c: \(\mathbf{I} \to \mathbf{S}_{1}\) 是道路 \(t \mapsto e^{\pi it}\)；\(\varphi(c)\) 的类是群 \(\pi_1(\mathbf{P}_{1}(\mathbf{R}), \varphi(1))\) 的一个生成元，该群同构于 \(\mathbf{Z}\)。

对每个整数 \(n \geqslant 2\) 及 \(S_{n}\) 的每一点 x，典范映射 \(\varphi: \mathbf{S}_{n} \to \mathbf{P}_{n}(\mathbf{R})\) 使 \((\mathbf{S}_{n}, x)\) 成为 \((\mathbf{P}_{n}(\mathbf{R}), \varphi(x))\) 的一个万有覆叠。对 \(S_{n}\) 中连接 x 与 -x 的每条道路 c，\(\varphi \circ c\) 的类生成群 \(\pi_{1}(\mathbf{P}_{n}(\mathbf{R}), \varphi(x))\)，该群同构于 Z/2Z。

\subitemhead{空间关于离散群逆紧作用的商}

引理 1. --- 设 X 是一个拓扑空间，G 是一个逆紧地作用于 X 上的离散群，p 是从 X 到 X/G 上的典范投影。

设 \(x\) 是 \(X\) 的一点，\(K_x\) 是它在 \(G\) 中的固定子。设 \(U_1\) 是 \(x\) 在 \(X\) 中的一个邻域；存在一个邻域 \(U\)（\(x\) 在 \(X\) 中），它在 \(K_x\) 下稳定且包含于 \(U_1\)，使 \(g \cdot U \cap U = \varnothing\)（对 \(g \notin K_x\)），且使典范映射 \(p\) 诱导 \(U / K_x\) 到邻域 \(V\)（\(p(x)\) 在 \(X / G\) 中的邻域）上的一个同胚。

由 TG, III, p. 32 的命题 8，存在一个邻域 \(U_{2}\)（x 在 X 中），它在 \(K_{x}\) 下稳定，使得 \(g \cdot U_{2} \cap U_{2} = \varnothing\)（对 \(g \notin K_{x}\)），且使得典范映射 p 诱导一个从 \(U_{2}/K_{x}\) 到 X/G 中 \(p(x)\) 的某个邻域上的同胚。

由于典范映射 \(U_{2} \rightarrow U_{2}/K_{x}\) 是闭的（TG, III, p. 28, 命题 2），存在 x 的一个开邻域 U，它包含于 \(U_{1} \cap U_{2}\) 且在 \(K_{x}\) 下稳定。\(U_{2}\) 中由 \(K_{x}\) 定义的等价关系还是开的（TG, III, p. 10, 引理 2）。于是由 TG, I, p. 32, 命题 4 推出，典范映射诱导一个从 \(U/K_{x}\) 到 X/G 中 \(p(x)\) 的某个开邻域上的同胚，引理得证。

命题 8. --- 设 X 是一个拓扑空间，G 是一个逆紧地作用于 X 上的离散群。用 p 表示从 X 到 X/G 的典范投影。

\begin{enumerate}
\def\labelenumi{\alph{enumi})}
\item
  假设 X 道路连通，且群 G 由 X 的各点的稳定子生成。那么典范同态 \(\pi_{1}(X,x)\to\pi_{1}(X/G,p(x))\) 对 X 的每一点 x 都是满射的。特别地，若 X 道路单连通，则 X/G 亦然。
\item
  若空间 X 是可解圈的，则空间 X/G 是可解圈的。
\item
  设 N 是由满足下列条件的元素 \(g \in G\) 组成的集合：存在一条道路 \(c_{g} \colon I \to X\)，使 \(c_{g}(0) = x\)、\(c_{g}(1) = g \cdot x\)，且使 X/G 中的圈 \(p \circ c_{g}\) 严格同伦于一条常值圈；N 是 G 的一个子群。设 \(g \in G\)，且存在 X 的一点 y 使 \(g \cdot y = y\)；取一条连接 x 与 y 的道路 c: I → X；于是道路 \(c' = c * (g \cdot c)\) 满足 \(c'(0) = x\)、\(c'(1) = g \cdot x\) 且 \([p \circ c'] = [p \circ c][p \circ (g \cdot c)]^{-1} = e_{p(x)}\)，因此 \(p \circ c'\) 严格同伦于一条常值圈。由于 G 由 X 的各点的固定子生成，由此得出 N = G。
\end{enumerate}

设 \(c\) 是 X/G 中 \(p(x)\) 处的一个圈。由 III, p. 287 的定理 4，存在一条道路 \(\widetilde{c} \colon \mathbf{I} \to \mathrm{X}\)，它提升 \(c\) 且使 \(\widetilde{c}(0) = x\)。由于 \(p(\widetilde{c}(1)) = c(1) = p(x)\)，存在 \(g \in \mathrm{G}\) 使 \(\widetilde{c}(1) = g \cdot x\)。选取一条道路 \(c_g \colon \mathbf{I} \to \mathrm{X}\)，它连接 \(x\) 与 \(g \cdot x\)，且使 \(p \circ c_g\) 严格同伦于一条常值圈。于是道路 \(c' = \widetilde{c} * \overline{c_g}\) 是 X 中 \(x\) 处的一个圈，且 \([p \circ c'] = [c]\)。这表明同态 \(\pi_1(\mathrm{X}, x) \to \pi_1(\mathrm{X/G}, p(x))\) 是满射的。另一断言立即得证。

现在来证明断言 \(b\)。空间 X/G 是局部道路连通的（III, p. 261, 命题 8）。

设 \(x\) 是 X 的一点，\(K_x\) 是它在 G 中的稳定子。设 \(U_1\) 是 \(x\) 的一个开邻域，使典范同态 \(\pi_1(U_1, x) \to \pi_1(X, x)\) 的像退缩为单位元。由引理 1（IV, p. 349），存在 X 中 \(x\) 的一个邻域 U，它包含于 \(U_1\)、在 \(K_x\) 下稳定，使 \(g \cdot U \cap U = \varnothing\)（对 \(g \notin K_x\)），且使典范映射 \(p\) 诱导 \(U/K_x\) 到 X/G 中 \(p(x)\) 的邻域 V 上的一个同胚。

设 \(c\) 是 V 中 \(p(x)\) 处的一个圈；把定理 4（III, p. 287）应用于拓扑空间 U 与群 \(K_x\)，存在一条道路 \(\widetilde{c} \colon \mathbf{I} \to \mathbf{U}\)，它使 \(\widetilde{c}(0) = x\) 且提升 \(c\)。必然有 \(\widetilde{c}(1) = x\)，因而 \(\widetilde{c}\) 是 U 中 \(x\) 处的一个圈。依假设，\(\widetilde{c}\) 在 X 中严格同伦于一条常值圈；于是 \(c\) 亦然，且典范同态 \(\pi_1(V, p(x)) \to \pi_1(X/G, p(x))\) 平凡。这表明：若 X 是可解圈的，则 X/G 是可解圈的。

\begin{enumerate}
\def\labelenumi{\arabic{enumi})}
\setcounter{enumi}{4}
\tightlist
\item
  在欧氏平面 \(R^{2}\) 中，设 P 为以 \((1/n,0)\) 为中心且过原点的诸圆（对 \(n \in N^{*}\)）之并构成的拓扑空间。空间 P 是紧的、连通的且局部道路连通的，但不是可解圈的。群 \(\pi_{1}(P,0)\) 赋以容许拓扑后是豪斯多夫的且非离散（III, p. 337, 习题 7）。
\end{enumerate}

注 1. --- III, p. 287 的定理 4 与命题 8（IV, p. 349）在更一般的假设下仍然成立：G 是 R 上的有限维李群且逆紧地作用于 X 上。更多细节参见 D. MONTGOMERY 与 C. T. YANG，“The existence of a slice”，Annals of mathematics 65 (1957), p. 108--116；R. PALAIS，“On the existence of slices for actions of non-compact Lie groups”，Annals of mathematics 73 (1961), p. 295--323；以及 G. BREDON，Introduction to compact transformation groups，Academic Press, 1972。

\section{可解圈空间的庞加莱群}

\subsection{\texorpdfstring{同态 \(\pi_{1}(f,a)\) 的性质}{同态 pi\_1(f,a) 的性质}}

设 A 与 B 是拓扑空间，\((\mathrm{E},p)\) 是 A 的一个覆叠，\((\mathrm{E}^{\prime},p^{\prime})\) 是 B 的一个覆叠，\(f\colon A\to B\) 与 \(g\colon E\to E^{\prime}\) 是满足 \(p^{\prime}\circ g=f\circ p\) 的连续映射。设 a 是 A 的一点，并设 \(b=f(a)\)。对所有 \(\gamma\in\pi_{1}(\mathrm{A},a)\) 及每个 \(x\in E_{a}\)，依 III, p. 304 的关系式 (3)，我们有 \(g(x\cdot\gamma)=g(x)\cdot f_{*}(\gamma)\)。若图表

\[
\begin{array}{c} \mathrm{E} \xrightarrow {g} \mathrm{E} ^ {\prime} \\ \Big \downarrow p \\ \mathrm{A} \xrightarrow {f} \mathrm{B} \end{array}
\]

是笛卡尔方块，则映射 g 诱导一个从 \(E_{a}\) 到 \(E'_{b}\) 上的双射。于是 \(\pi_{1}(A,a)\) 在纤维 \(E_{a}\) 上的作用，是群 \(\pi_{1}(B,b)\) 在 \(E'_{b}\) 上的作用与同态 \(\pi_{1}(f,a)\)（从 \(\pi_{1}(A,a)\) 到 \(\pi_{1}(B,b)\) 中）的复合。

命题 1. --- 设 A 是可解圈拓扑空间，B 是局部道路连通拓扑空间，并设 \(f\colon\mathrm{A}\to\mathrm{B}\) 是连续映射。设 a 是 A 的一点，\(b=f(a)\)。假设 A 的每个覆叠都同构于 B 的某个覆叠在 \(f\) 下的逆像。那么同态 \(\pi_{1}(f,a)\colon\pi_{1}(\mathrm{A},a)\to\pi_{1}(\mathrm{B},b)\) 是单射的。

由于空间 A 是可解圈的，存在 A 的一个覆叠 E，其纤维 \(E_{a}\) 是主齐性 \(\pi_{1}(A,a)\) -集合（IV, p. 342, 定理 1）。依假设，这个作用经由同态 \(\pi_{1}(f,a)\) 由 \(\pi_{1}(B,b)\) 在 \(E_{a}\) 上的一个作用诱导。这蕴含同态 \(\pi_{1}(f,a)\) 的单射性。

命题 2. --- 设 A 是连通且局部道路连通的拓扑空间，B 是可解圈拓扑空间，并设 \(f\colon\mathrm{A}\to\mathrm{B}\) 是连续映射。设 a 是 A 的一点，\(b=f(a)\)。下列断言等价：

\begin{enumerate}
\def\labelenumi{(\roman{enumi})}
\item
  同态 \(\pi_1(f,a)\colon \pi_1(A,a)\to \pi_1(B,b)\) 是满射的。
\item
  对 B 的覆叠的每个偶 (E, E')，映射
\end{enumerate}

\[
f ^ {*} \colon \mathscr {C} _ {\mathrm{B}} (\mathrm{E}; \mathrm{E} ^ {\prime}) \rightarrow \mathscr {C} _ {\mathrm{A}} (\mathrm{A} \times_ {\mathrm{B}} \mathrm{E}; \mathrm{A} \times_ {\mathrm{B}} \mathrm{E} ^ {\prime})
\]

它把 B-态射 \(g\colon \mathrm{E}\to \mathrm{E}^{\prime}\) 对应到 A-态射

\[
f ^ {*} (g) \colon \mathrm{A} \times_ {\mathrm{B}} \mathrm{E} \rightarrow \mathrm{A} \times_ {\mathrm{B}} \mathrm{E} ^ {\prime}, \quad (x, y) \mapsto (x, g (y))
\]

是双射的。

借助同态 \(\pi_1(f, a)\)，每个 \(\pi_1(B, b)\) -集合都赋有一个 \(\pi_1(A, a)\) -集合的结构。于是条件 (i) 等价于下列条件 (i')：

(i') 对每个偶 \((\mathrm{F},\mathrm{F}^{\prime})\)（\(\pi_{1}(\mathrm{B},b)\) -集合的偶），每个 \(\pi_{1}(\mathrm{A},a)\) -态射（从 F 到 \(F^{\prime}\)）都是 \(\pi_{1}(\mathrm{B},b)\) -态射。事实上，若同态 \(\pi_{1}(f,a)\) 是满射的，则每个 \(\pi_{1}(\mathrm{A},a)\) -态射（\(\pi_{1}(\mathrm{B},b)\) -集合的态射）都是 \(\pi_{1}(\mathrm{B},b)\) -态射。反之，取一个退缩为一点的 \(\pi_{1}(\mathrm{B},b)\) -集合 F，并置 \(\mathrm{F}^{\prime}=\pi_{1}(\mathrm{B},b)/f_{*}(\pi_{1}(\mathrm{A},a))\)。从 F 到 \(F^{\prime}\) 的、以 \(f_{*}(\pi_{1}(\mathrm{A},a))\) 为像的映射是一个 \(\pi_{1}(\mathrm{A},a)\) -态射，但不是 \(\pi_{1}(\mathrm{B},b)\) -态射，如果 \(F^{\prime}\) 不退缩为一点，也就是说，如果 \(\pi_{1}(f,a)\) 不是满射的。

由于空间 B 是可解圈的，每个 \(\pi_{1}(B,b)\) -集合都同构于 \(\pi_{1}(B,b)\) -集合 \(E_{b}\)，其中 E 是 B 的覆叠（IV, p. 344, 注 1）。于是 (i') 与 (ii) 的等价性由 III, p. 310 的命题 2 得出。

命题 3. --- 设 B 是可解圈拓扑空间，A 是 B 的连通且局部道路连通子空间（例如一个连通开子集）。设 a 是 A 的一点。下列性质等价：

\begin{enumerate}
\def\labelenumi{(\roman{enumi})}
\item
  B 的每个覆叠都在 A 上可平凡化。
\item
  由典范单射导出的、从 \(\pi_{1}(\mathrm{A},a)\) 到 \(\pi_{1}(\mathrm{B},a)\) 中的同态的像退缩为单位元。
\end{enumerate}

蕴涵关系 (ii)⇒(i) 由推论 3（III, p. 309）得出。

反之，假设 B 的每个覆叠都在 A 上可平凡化。对道路单连通覆叠 \((\lambda_{a}(\mathrm{B}),\varepsilon_{\mathrm{B}})\)（IV, p. 342, 定理 1）而言尤其如此，从而同态 \(\pi_{1}(\mathrm{A},a)\to\pi_{1}(\mathrm{B},a)\) 平凡（III, p. 309, 命题 1 的推论 3）。

\subsection{相对连通映射}

定义 1. --- 设 X 与 Y 是拓扑空间，f 是从 X 到 Y 的连续映射。称映射 f 为相对连通的，如果 Y 的每一点都有一个由满足下列条件的集合 V 组成的基本邻域系：\(f^{-1}(V)\) 连通且非空。

设 \(f \colon X \to Y\) 是连续映射；为使 \(f\) 相对连通，必要且充分的是 \(Y\) 的每一点都有一个邻域 \(V\)，使映射 \(f_V \colon f^{-1}(V) \to V\)（由 \(f\) 导出）相对连通。

设 X 与 Y 是拓扑空间，\(f: X \to Y\) 是相对连通连续映射。f 的像在 Y 中稠密。对 Y 的每个开子集 V，映射 \(f_{\mathrm{V}}: f^{-1}(\mathrm{V}) \to \mathrm{V}\) 相对连通。

命题 4. --- 设 X 与 Y 是拓扑空间，\(f: X \to Y\) 是相对连通连续映射。

\begin{enumerate}
\def\labelenumi{\alph{enumi})}
\item
  对 X 的每个连通分支 U，\(f(\overline{\mathrm{U}})\) 是 Y 的连通、开且闭的子集，且我们有 \(f^{-1}(\overline{f(\mathrm{U})}) = \mathrm{U}\)。
\item
  对 Y 的每个连通分支 V，存在 X 的一个连通分支 U，使 \(V = \overline{f(U)}\)。由 f 取子空间导出的、从 U 到 V 的映射相对连通。
\item
  X（相应地，Y）的连通分支都是开且闭的。
\item
设 V 是由满足下列条件的点 \(y \in Y\) 组成的集合：它们有一个邻域 W，使 \(f^{-1}(W)\) 连通且与 U 相交；V 是 Y 的开子集。它包含 \(\overline{f(U)}\)，因为 f 相对连通。反之，设 \(y \in V\)；设 W 是 \(y\) 的一个邻域，使 \(f^{-1}(W)\) 连通且与 U 相交。对每个邻域 \(W'\)（\(y\) 的、使 \(f^{-1}(W')\) 连通的邻域），\(f^{-1}(W \cap W')\) 非空，因此 \(f^{-1}(W \cup W')\) 连通（TG, I, p. 81, 命题 2）。依假设，\(f^{-1}(W \cup W')\) 与 U 相交；于是 \(f^{-1}(W \cup W') \subset U\)，且更有 \(W' \cap f(U) \neq \varnothing\)。这证明了 \(y \in \overline{f(U)}\)；因此 \(V = \overline{f(U)}\)。特别地，集合 \(\overline{f(U)}\) 在 Y 中既开且闭；此外，命题 1（TG, I, p. 81）蕴含它连通。
\end{enumerate}

前面的论证还表明 \(f^{-1}(\overline{f(\mathrm{U})}) \subset \mathrm{U}\)，从而有等式 \(f^{-1}(\overline{f(\mathrm{U})}) = \mathrm{U}\)，另一个包含关系是显然的。特别地，U 在 X 中既开且闭。

\begin{enumerate}
\def\labelenumi{\alph{enumi})}
\setcounter{enumi}{1}
\tightlist
\item
设 V 是 Y 的一点 y 的连通分支，W 是 y 的一个使 \(f^{-1}(W)\) 连通非空的邻域，U 是包含 \(f^{-1}(W)\) 的 X 的连通分支。由于 \(y \in \overline{f(U)}\)，由 a) 及 y 的连通分支的定义得出 \(\mathrm{V} = \overline{f(\mathrm{U})}\)。因此 V 在 Y 中既开且闭。
\end{enumerate}

推论 1. --- 设 X 与 Y 是拓扑空间，\(f: X \to Y\) 是连续且相对连通的映射。通过取连通分支，映射 \(f\) 诱导一个从 X 的连通分支所成之集到 Y 的连通分支所成之集上的双射。

推论 2. --- 设 X 与 Y 是拓扑空间，\(f: X \to Y\) 是连续映射。为使 \(f\) 相对连通，必要且充分的是下列三条性质成立：

\begin{enumerate}
\def\labelenumi{\alph{enumi})}
\item
  像 f(X) 在 Y 中稠密；
\item
  空间 Y 局部连通；
\item
  对 Y 的每个开连通集合 V，集合 \(f(\mathrm{V})\) 连通。
\end{enumerate}

假设映射 f 相对连通。\(f(\mathrm{X})\) 在 Y 中的稠密性由相对连通映射的定义得出。设 V 是 Y 的开子集；映射 \(f_{\mathrm{V}}: f^{-1}(\mathrm{V}) \to \mathrm{V}\) 相对连通。由命题 4，V 的连通分支在 V 中既开且闭。由此得出 Y 局部连通（TG, I, p. 85, 命题 11）。为证明断言 c)，只需证明：若 Y 连通，则 X 连通。由引理，存在 X 的一个连通分支 U，使 \(Y = \overline{f(U)}\) 且 \(f^{-1}(\overline{f(U)}) = U\)，于是 U = X，故 X 连通。

反之，假设条件 \(a\)、\(b\)、\(c\) 满足，我们来证明 f 相对连通。设 \(y\) 是 Y 的一点；由于 Y 局部连通，\(y\) 有一个由连通开邻域组成的基本邻域系。若 W 是这样一个邻域，则条件 \(c\) 与 \(a\) 蕴含 \(f^{-1}(W)\) 连通且非空。因此映射 f 相对连通。

推论 3. --- 设 X 与 Y 是拓扑空间，\(f\colon X\to Y\) 是连续且相对连通的映射。设 F 是一个集合，\(g\colon X\to F\) 是局部常值映射。存在唯一的局部常值映射 \(h\colon Y\to F\)，使 \(g=h\circ f\)。

g 在 X 的每个连通分支上的限制是常值的。由推论 1，存在一个映射 \(h: Y \to F\)，它在 Y 的每个连通分支上取常值，使 \(g = h \circ f\)。映射 h 是局部常值的，因为 Y 的连通分支都是开的。这样一个映射的唯一性由 \(f(X)\) 在 Y 中稠密这一事实得出。

例. --- 1) 设 X 是拓扑空间，R 是 X 上的一个等价关系。记 Y 为商拓扑空间 X/R，记 \(f\colon X \to Y\) 为典范映射。假设 R 的等价类都连通。那么，对 Y 的每个开连通子集 V，集合 \(f^{-1}(V)\) 连通（TG, I, p. 23, 推论 1 及 p. 82, 命题 7）。若空间 Y 局部连通，则映射 f 因而相对连通。

\begin{enumerate}
\def\labelenumi{\arabic{enumi})}
\setcounter{enumi}{1}
\tightlist
\item
  设 Y 是局部道路连通拓扑空间，X 是 Y 的开子集。X 中道路的空间 \(\mathcal{C}_{c}(\mathbf{I};\mathbf{X})\) 与 Y 中道路的空间 \(\mathcal{C}_{c}(\mathbf{I};\mathbf{Y})\) 的一个子空间等同起来；假设它在其中稠密。于是 X 到 Y 中的典范单射是相对连通映射。
\end{enumerate}

事实上，只需证明：对 Y 的每个非空连通开子集 V，集合 \(V \cap X\) 连通且非空。空间 \(\mathcal{C}_{c}(\mathbf{I};\mathrm{V})\) 是 \(\mathcal{C}_{c}(\mathbf{I};\mathrm{Y})\) 的一个非空开集；它与 \(\mathcal{C}_{c}(\mathbf{I};\mathrm{X})\) 相交，这就证明了 \(V \cap X \neq \varnothing\)。设 x 与 \(x'\) 是 \(V \cap X\) 的点。由于 \(V \cap X\) 局部道路连通，点 x 与 \(x'\) 有开邻域 U 与 \(U'\)，它们道路连通且包含于 \(V \cap X\)。由于 V 道路连通（III, p. 260, 命题 7），存在 V 中连接 x 与 \(x'\) 的道路。由稠密性假设及紧收敛拓扑的定义，存在 \(V \cap X\) 中连接 U 的一点与 \(U'\) 的一点的一条道路。于是存在一条连接 x 与 \(x'\) 的、\(V \cap X\) 中的道路，这就证明了集合 \(V \cap X\) 连通。

命题 5. --- 设 X 与 Y 是拓扑空间，\(f: X \to Y\) 是连续且相对连通的映射。对 Y 的覆叠的每个偶 \((T, T')\)，映射 \(f^{*}: \mathcal{C}_{Y}(T; T') \to \mathcal{C}_{X}(X \times_{Y} T; X \times_{Y} T')\) 是双射的。

设 \(\mathcal{F}\) 是 X 上由从 \(X \times_{Y} T\) 到 \(X \times_{Y} T'\) 中的 X-态射构成的层，\(\mathcal{G}\) 是 Y 上由从 T 到 \(T'\) 中的 Y-态射构成的层（I, p. 45, 例 4）。对 Y 的每个开集 U，置 \(\varphi_{\mathrm{U}} = (f_{\mathrm{U}})^{*}: \mathcal{G}(\mathrm{U}) \to \mathcal{F}(f^{-1}(\mathrm{U}))\)。诸映射 \(\varphi_{\mathrm{U}}\) 定义一个层的态射 \(\varphi: \mathcal{G} \to \varphi_{*}(\mathcal{F})\)，只需证明 \(\varphi\) 是层的一个同构。

由于 Y 局部连通（IV, p. 354, 推论 2），T 与 T' 在其上可平凡化的连通开集构成 Y 的拓扑的一个基。由 I, p. 55 的推论 2，只需证明：对这样一个开集 U，映射 \(\varphi_{\mathrm{U}}\) 是双射的，这使我们可以假设 Y 连通，且覆叠 T 与 T' 是平凡覆叠 \(Y \times F\) 与 \(Y \times F'\)，其中 F 与 F' 是赋有离散拓扑的集合。映射 \((x, (y, t)) \mapsto (x, t)\) 把 X-空间 \(X \times_{Y} (Y \times F)\) 与 \(X \times F\) 等同起来（相应地，把 X-空间 \(X \times_{Y} (Y \times F')\) 与 \(X \times F'\) 等同起来）。由于空间 X 连通（IV, p. 354, 推论 2），集合 \(\mathcal{C}_{\mathrm{Y}}(Y \times F; Y \times F')\) 与 \(\mathcal{C}_{\mathrm{X}}(X \times F; X \times F')\) 二者都与从 F 到 F' 中的映射所成的集合 \(\mathcal{F}(F;F')\) 等同，且映射 \(f^{*}\) 与 \(\mathcal{F}(F;F')\) 的恒等映射等同。证明到此完成。

推论 1. --- 设 X 与 Y 是拓扑空间，\(f: X \to Y\) 是相对连通连续映射。设 T 与 T' 是 Y 的覆叠。若 X 的覆叠 \(X \times_Y T\) 与 \(X \times_Y T'\) 同构，则覆叠 T 与 T' 同构。

事实上，设 \(h \colon X \times_{Y} T \to X \times_{Y} T'\) 与 \(h' \colon X \times_{Y} T' \to X \times_{Y} T\) 是互为逆的 X-同构。由命题 5，存在 Y-态射 \(g \colon \mathrm{T} \to \mathrm{T}'\) 与 \(g' \colon \mathrm{T}' \to \mathrm{T}\)，使 \(f^{*}(g) = h\) 且 \(f^{*}(g') = h'\)。于是有 \(f^{*}(g' \circ g) = f^{*}(g') \circ f^{*}(g) = \mathrm{Id}_{\mathrm{X} \times_{\mathrm{Y}} \mathrm{T}}\)，因此 \(g' \circ g = \mathrm{Id}_{\mathrm{T}}\)，因为映射 \(f^{*}\) 是单射的。类似地，\(g \circ g' = \mathrm{Id}_{\mathrm{T}'}\)。因此覆叠 T 与 \(\mathrm{T}'\) 同构。

推论 2. --- 设 X 与 Y 是拓扑空间，\(f: X \to Y\) 是连续且相对连通的映射。若空间 X 单连通，则空间 Y 亦然。

事实上，由推论 1，Y 的每个覆叠都可平凡化。

推论 3. --- 设 X 是局部道路连通拓扑空间，Y 是可解圈拓扑空间，\(f: X \to Y\) 是相对连通连续映射。对 X 的每一点 \(x\)，同态 \(\pi_1(f, x): \pi_1(X, x) \to \pi_1(Y, f(x))\) 是满射的。

由 IV, p. 353 的命题 4，我们可以假设空间 X 与 Y 都连通。于是推论由 IV, p. 352 的命题 5 与命题 2 得出。

命题 6. --- 设 X 与 Y 是拓扑空间，\(f: \mathrm{X} \to \mathrm{Y}\) 是连续且相对连通的映射。假设 Y 的每一点都有一个开邻域 V，其逆像 \(f^{-1}(V)\) 单连通。那么，对 X 的每个覆叠 Z，存在 Y 的一个覆叠 T，使 \(X \times_{Y} T\) X-同构于 Z。

设 \(\mathcal{U}\) 是 Y 中其逆像在 X 中单连通的开子集所成的集合。对每个 \(V \in \mathcal{U}\)，覆叠 \(f^{-1}(V) \times_{X} Z\)（\(f^{-1}(V)\) 的覆叠）是可平凡化的；于是存在一个离散空间 \(F_V\) 及覆叠的一个同构 \(g_V: f^{-1}(V) \times_{X} Z \to f^{-1}(V) \times F_V\)。对开集的每个偶 \((V, V')\)（属于 \(\mathcal{U}\)），映射 \(f_{V \cap V'}: f^{-1}(V \cap V') \to V \cap V'\) 相对连通。由 IV, p. 356 的命题 5，存在 \(V \cap V'\) 的覆叠的唯一同构 \(h_{V',V}: (V \cap V') \times F_V \to (V \cap V') \times F_{V'}\)，使得有 \(f^*(h_{V',V})(x,t) = g_{V'}(g_V^{-1}(x,t))\)（对所有 \(x \in V \cap V'\) 及每个 \(t \in F_V\)）。若 V、\(V'\)、\(V''\) 是 \(\mathcal{U}\) 的元素，则有 \(h_{V'',V}(x,t) = h_{V'',V'}(h_{V',V}(x,t))\)（对所有 \(x \in V \cap V' \cap V''\) 及每个 \(t \in F_V\)）。于是存在唯一的 Y-空间 T，且对每个 \(V \in \mathcal{U}\)，存在同构 \(h_V: T_V \to V \times F_V\)，使得有 \(h_{V',V}(x,t) = h_{V'} \circ h_V^{-1}(x,t)\)——其中偶 \((V, V')\) 属于 \(\mathcal{U}\)、\(x \in V \cap V'\) 任意、\(t \in F_V\) 任意（参见 TG, I, p. 16）。

空间 T 特别地是 Y 的一个覆叠。此外，存在唯一的、到 Z 上的 \(X \times_{Y} T\) 的映射，它在 \(f^{-1}(V) \times_{Y} T\) 上的限制由 \(g_{V}^{-1} \circ f^{*}(h_{V})\) 给出，且它是 X-空间的一个同构，命题由此得证。

推论. --- 设 X 与 Y 是可解圈拓扑空间，\(f: X \to Y\) 是连续且相对连通的映射。假设 Y 的每一点都有一个邻域 V，其逆像 \(f^{-1}(V)\) 单连通。那么，对 X 的每一点 x，同态 \(\pi_{1}(f,x): \pi_{1}(X,x) \to \pi_{1}(Y,f(x))\) 是双射的。

可以假设空间 X 与 Y 都连通（IV, p. 353, 命题 4）。由推论 3（IV, p. 357），同态 \(\pi_1(f,x)\) 是满射的。命题 6 与 IV, p. 351 的命题 1 蕴含它是单射的。

注. --- 设 Y 是一个拓扑空间，X、X' 与 Y' 是 Y 的子空间，满足 X ⊂ X' ⊂ Y' ⊂ Y。假设 X 到 Y 中的典范单射相对连通。对 Y 的每个连通开子集 V，集合 V ∩ X 连通且在 V 中稠密（IV, p. 354, 推论 2）；于是集合 V ∩ X' 连通（TG, I, p. 81, 命题 1）。这证明了 X' 到 Y' 中的典范单射相对连通。

设 V 是 Y 的一个开子集；由前文，\(V \cap X\) 到 \(V \cap X'\) 中的典范单射相对连通。若集合 \(V \cap X\) 单连通，则 \(V \cap X'\) 亦然，这由 IV, p. 357 的推论 2 得出。因此，若 X 到 Y 中的典范单射满足命题 6 的假设，则 \(X'\) 到 \(Y'\) 中的典范单射亦然。

例. --- 设 Y 是一个局部有限维微分流形，Z 是 Y 的一个闭子流形（VAR, R, 5.8.3）。置 X = Y - Z。

\begin{enumerate}
\def\labelenumi{\alph{enumi})}
\item
  若 Z 在每点的余维数至少为 2，则 X 到 Y 中的典范单射相对连通。事实上，设 z 是 Z 的一点；存在 z 在 Y 中的开邻域 V 及 V 到 R 上的一个有限维向量空间 E 上的同胚 \(\varphi\)，使 \(\varphi(V \cap Z)\) 是 E 的一个余维数 \(\geqslant 2\) 的向量子空间 F。集合 E-F 连通（TG, VI, p. 5, 命题 4），断言由此得证。
\item
  进而假设 Z 在每点的余维数至少为 3；于是命题 6 及其推论的假设满足，因为沿用前段的记号，E-F 单连通（I, p. 128, 例 4）。
\end{enumerate}

由于微分流形是可解圈空间（IV, p. 347），本节的结果有如下的特例。

设 Y 是一个局部有限维微分流形，Z 是 Y 的一个闭子流形，i 是 Y−Z 到 Y 中的典范单射。

\begin{enumerate}
\def\labelenumi{\alph{enumi})}
\item
  若 Z 在 Y 中的余维数在每点 ≥ 1，则流形 Y−Z 在 Y 中稠密，且映射 \(\pi_{0}(i)\) 是满射的。
\item
  若 Z 在 Y 中的余维数在每点 ≥ 2，则映射 \(\pi_{0}(i)\) 是双射的，且对 Y-Z 的每一点 x，映射 \(\pi_{1}(i,x)\) 是满射的。
\item
  若 Z 在 Y 中的余维数在每点 \(\geqslant 3\)，则对 Y-Z 的每一点 x，映射 \(\pi_{0}(i)\) 与 \(\pi_{1}(i,x)\) 都是双射的。
\end{enumerate}

\subsection{庞加莱群的表现}

定理 1. --- 设 X 是一个紧的可解圈拓扑空间，x 是 X 的一点。庞加莱群 \(\pi_{1}(X,x)\) 是有限表现的。

x 在 X 中的弧连通分支是开的、闭的且可解圈的；这使我们可以假设空间 X 连通。由于空间 X 可解圈，X-空间 \(E = \lambda_{x}(X)\)——赋以终点映射——是一个非空且道路单连通的覆叠（IV, p. 342, 定理 1）。群 \(G = \text{Aut}_{\text{X}}(\text{E})\) 同构于 \(\pi_{1}(\text{X}, x)\)；因此只需证明群 G 是有限表现的。

由于 X 紧，X 的每一点 x 都有一个紧邻域 \(K_{x}\)，覆叠 E 在其上可平凡化（TG, I, p. 65, 推论）。由于 X 局部连通，每个 \(x \in X\) 都有一个开邻域 \(W_{x}\)（连通，且包含于 \(K_{x}\)）。设 F 是 X 的一个有限子集，使得诸 \(W_{x}\)（\(x \in F\)）覆盖 X。设 n 为 F 的基数。

我们用归纳法证明：对每个整数 \(k\)（\(1 \leqslant k \leqslant n\)），存在一个基数为 \(k\)、包含于 F 的子集 A，且对 \(x \in A\)，存在一个截面 \(s_x\)（\(p\) 在 \(K_x\) 上的截面），使得诸 \(s_x(W_x)\)（对 \(x \in A\)）之并是 E 的一个连通子集。断言对 \(k = 1\) 成立。假设它对一个整数 \(k\)（\(1 \leqslant k < n\)）成立，我们来证明它对 \(k + 1\) 成立。设 A 是 F 的基数为 \(k\) 的子集，且对每个 \(x \in A\)，设 \(s_x\) 是 E 在 \(K_x\) 上的一个截面，使 \(\bigcup_{x \in A} s_x(W_x)\) 连通。开集 \(\bigcup_{x \in A} W_x\) 与 \(\bigcup_{x \in F-A} W_x\) 非空且覆盖 X；由于 X 连通，它们的交非空。于是存在 \(x \in A\)、\(y \in F - A\) 及 \(z \in W_x \cap W_y\)。置 \(A' = A \cup \{y\}\)，并选取一个截面 \(s_y\)（E 在 \(K_y\) 上的截面），使 \(s_y(z) = s_x(z)\)。连通开集 \(\bigcup_{p \in A} s_p(W_p)\) 与 \(s_y(W_y)\) 非空且有公共点；因此其并连通。这就对 \(k + 1\) 证明了断言。由归纳法，它于是对每个整数 \(k \in \{1, \ldots, n\}\) 成立。

把上述结果应用于 k = n，于是对每个 \(x \in F\)，存在一个截面 \(s_x\)（E 在 \(K_x\) 上的截面），使得 \(U = \bigcup_{x \in F} s_x(W_x)\) 是 E 的一个开连通子集。我们有 \(p(U) = X\)。由于群 G 在覆叠 E 的每条纤维中传递地作用，我们有 GU = E。

U 的闭包包含于 \(\bigcup_{x\in\mathrm{F}}s_x(\mathrm{K}_x)\)，因而是紧的。G 在 E 上的作用是逆紧的（I, p. 96, 推论 1），故由偶 \((g,x)\in\mathrm{G}\times\mathrm{E}\)（满足 \(x\in\overline{\mathrm{U}}\) 且 \(gx\in\overline{\mathrm{U}}\)）所成之集是 \(\mathrm{G}\times\mathrm{E}\) 的一个紧子集（TG, I, p. 77, 命题 6）。由此得出，由 \(g\in\mathrm{G}\)（满足 \(\overline{\mathrm{U}}\cap g\overline{\mathrm{U}}\neq\varnothing\)）所成之集是紧的。由于它离散，它有限。更有，由 \(g\in\mathrm{G}\)（满足 \(\mathrm{U}\cap g\mathrm{U}\neq\varnothing\)）所成之集有限。定理于是由 I, p. 136 的命题 10 得出。

\subsection{关于波兰空间的补充}

引理 1. --- 设 X 是一个拓扑空间，A 是 X 的一个子集。为使 A 是贫集，必要且充分的是：存在 X 的一个开覆盖 \((\mathrm{U}_{i})_{i\in\mathrm{I}}\)，使 \(A\cap U_{i}\) 在 \(U_{i}\) 中贫（对每个 \(i\in I\)）。

设 \(\mathcal{O}\) 是 X 中使 \(A \cap U\) 贫的开子集 U 所成的集合。由两两不交的开集组成的 \(\mathcal{O}\) 的子集所成的族，按包含关系是归纳的。设 \(\mathfrak{U}\) 是一个极大元；这样一个元素的存在性由 E, III, p. 20, 定理 2 得出。

设 O 为 X 中属于 \(\mathfrak{U}\) 的诸开集之并。对每个开集 \(U \in \mathfrak{U}\)，设 \((B_{U,n})\) 是 U 的无处稠密子集的一个序列，其并为 \(A \cap U\)。对每个整数 n，设 \(B_n = \bigcup_{U \in \mathfrak{U}} B_{U,n}\)。

当 U 取遍诸 U 时，它相对于 O 是无处稠密的（TG, IX, p. 52, 命题 1）。由 TG, IX, p. 53, 命题 2，\(B_{n}\) 是 X 的一个无处稠密子集，因为 O 在 X 中开。于是，等于诸 \(B_{n}\) 之并的集合 A ∩ O 是 X 的一个贫子集。

设 F 为 O 的补集。它是 X 的闭子集；我们来证明它的内部为空。若不然，设 x 是 F 的一个内点。依假设，存在 x 的一个开邻域 V——可设它包含于 F——使 A ∩ V 在 X 中贫。于是 V 与 X 中属于 U 的诸开集不交，且 U ∪ \{V\} 是由属于 O 的两两不交的开集组成的一个集合，这与 U 的极大性矛盾。关系式

\[
\mathrm{A} = (\mathrm{A} \cap \mathrm{F}) \cup (\mathrm{A} \cap \mathrm{O})
\]

于是蕴含 A 贫，此即所要证明的。

设 X 是一个拓扑空间。

回忆（TG, IX, p. 69）：称 X 的一个子集 A 是可达的，如果存在 X 的一个开集 U，使 \(U \cap \complement A\) 与 \(A \cap \complement U\) 都在 X 中贫。X 的可达子集的全体是一个 \(\sigma\)-代数，它包含博雷尔 \(\sigma\)-代数（TG, IX, p. 69, 引理 8 及其证明）。贫子集是可达的。

对 X 的任一子集 A，设 D(A) 为由满足下列条件的点 \(x \in X\) 组成的集合：对 x 的每个邻域 V，A∩V 都不贫。我们还置 \(\mathrm{D}^{*}(\mathrm{A}) = \mathrm{A} \cup \mathrm{D}(\mathrm{A})\)。

引理 2. --- 设 X 是一个拓扑空间，A 是 X 的一个子集。

\begin{enumerate}
\def\labelenumi{\alph{enumi})}
\item
  集合 D(A) 在 X 中闭；其补集是 X 中使 A ∩ U 为贫集的最大开集 U。
\item
  为使 A 贫，必要且充分的是 D(A) 为空。
\item
  集合 \(\mathrm{D}^{*}(\mathrm{A})\) 是可达的。
\item
  对 X 的每个包含 A 的可达集 B，\(D^{*}(A) \cap \complement B\) 贫。
\end{enumerate}

设 U 为 X 中使 \(A \cap V\) 贫的诸开集 V 之并。为使一点 x 属于 U，必要且充分的是它有一个邻域 V 使 \(A \cap V\) 贫。于是 \(U = \complement D(A)\)，因此 D(A) 是 X 的闭子集。按构造，子空间 U 有一个由与 A 的交为贫的开集组成的覆盖。把引理 1 应用于拓扑空间 U 与子集 \(A \cap U\)，可知 \(A \cap U\) 在 U 中贫，因而也在 X 中贫。这就证明了断言 a)，因为按构造，X 中使 \(A \cap V\) 贫的每个开集 V 都包含于 U。

断言 b) 由此立得。

\begin{enumerate}
\def\labelenumi{\alph{enumi})}
\setcounter{enumi}{2}
\item
  我们有 \(D^{*}(A) = A \cup D(A) = (A \cap U) \cup D(A)\)。集合 \(D(A)\) 是闭的，因而是可达的（IV, p. 361）；贫的部分 \(A \cap U\) 也是可达的。于是 \(D^{*}(A)\) 可达（同前）。
\item
  设 B 是 X 的一个包含 A 的可达子集。其补集 \(\complement B\) 于是是 X 的可达子集（同前），因而存在 X 的一个开集 V，使 \(V \cap B\) 与 \(\complement V \cap \complement B\) 贫。由于 \(A \subset B\)，\(V \cap A\) 仍贫，故 \(V \subset U\)。由于 \(\complement U \cap \complement B\) 包含于 \(\complement V \cap \complement B\)，它也是 X 的贫子集。最后，包含关系
\end{enumerate}

\[
\mathrm{D} ^ {*} (\mathrm{A}) \cap \complement \mathrm{B} = ((\mathrm{A} \cap \mathrm{U}) \cap \complement \mathrm{B}) \cup (\complement \mathrm{U} \cap \complement \mathrm{B}) \subset (\mathrm{A} \cap \mathrm{U}) \cup (\complement \mathrm{U} \cap \complement \mathrm{B})
\]

证明了 \(D^{*}(A) \cap \complement B\) 是 X 的一个贫子集。

注 1. --- 设 G 是一个拓扑群，A 是 G 的一个贫子集。假设 A 包含 G 的一个非空开集 U，并设 y 是 U 的一点。那么 U 贫，且对每个 \(x \in G\)，\(xy^{-1}U\) 是 x 在 G 中的邻域。于是 G 的每一点都有一个贫邻域，故 G 贫（IV, p. 360, 引理 1）。

反之，若 G 不贫，则每个贫子集的内部为空，且 G 是贝尔空间。由于贝尔空间不贫，这就证明了下列断言：为使一个拓扑群是贝尔空间，必要且充分的是它不贫。

命题 7. --- 设 X 是一个豪斯多夫空间。X 的每个苏斯林子空间（TG, IX, p. 59, 定义 2）都是可达的。

设 S 是 X 的一个苏斯林子空间。按定义，存在一个具有可数基的完备度量空间 P 及一个从 P 到 S 上的连续满射 g。由 TG, IX, p. 64, 引理 3，存在度量空间 P 的一个筛 \(\mathrm{C} = (\mathrm{C}_{n}, p_{n}, \varphi_{n})_{n \in \mathbf{N}}\)。

对每个整数 \(n\) 及每个 \(c \in \mathrm{C}_n\)，记 \(\mathrm{F}_n(c)\) 为集合 \(\varphi_n(c)\) 在 \(g\) 下的像；再置 \(\mathrm{F}_n^*(c) = \mathrm{D}^*(\mathrm{F}_n(c))\)，以及

\[
\mathrm{G} _ {n} (c) = \mathrm{F} _ {n} ^ {*} (c) \cap \mathbb {C} \left(\bigcup_ {c ^ {\prime} \in p _ {n} ^ {- 1} (c)} \mathrm{F} _ {n + 1} ^ {*} (c ^ {\prime})\right).\tag{1}
\]

对每个 \(c \in C_{n}\)，\(\mathrm{F}_{n}^{*}(c)\) 是可达的（IV, p. 361, 引理 2, c)）。由于 \(C_{n+1}\) 可数，诸可达子集 \(\mathrm{F}_{n+1}^{*}(c')\)（\(c'\) 取遍 \(p_{n}^{-1}(c)\)）之并仍是 X 的一个可达子集。它包含诸 \(\mathrm{F}_{n+1}(c')\) 之并，后者等于 \(\mathrm{F}_{n}(c)\)。于是（同前， d)），\(\mathrm{G}_{n}(c)\) 是 X 的一个贫子集。因此，诸集合 \(\mathrm{G}_{n}(c)\)（对 \(n \in \mathbf{N}\) 及 \(c \in C_{n}\)）之并 G 是 X 的一个贫子集。

设 \(c_0 \in \mathrm{C}_0\)，设 \(x\) 是 \(\mathrm{F}_0^*(c_0) \cap \complement \mathrm{G}\) 的一个元素；我们来证明 \(x \in \mathrm{F}_0(c_0)\)。由于 \(x \notin \mathrm{G}_0(c_0)\) 且 \(x \in \mathrm{F}_0^*(c_0)\)，由关系式 (1)，存在 \(c_1 \in p_0^{-1}(c_0)\) 使 \(x \in \mathrm{F}_1^*(c_1)\)。由归纳法，存在一个元素 \(c = (c_n)_n\)（属于 \(\prod_n C_n\)），使得对每个整数 \(n \in \mathbf{N}\)，有 \(x \in \mathrm{F}_n^*(c_n)\) 且 \(p_n(c_{n+1}) = c_n\)。

诸子集 \(\varphi_n(c_n)\) 构成 P 中一个柯西滤子的基；由于 P 完备，这个滤子收敛于 P 的一点 \(p\)。这个滤子在 \(g\) 下的像是 X 上的一个滤子 F，以诸 \(\mathrm{F}_n(c_n)\)（\(n\in \mathbf{N}\)）组成的集合为基，且收敛于 \(g(p)\)。由于 \(\mathrm{F}_n^* (c_n)\) 包含于 \(\overline{\mathrm{F}_n(c_n)}\) 且 \(x\) 属于每个 \(\mathrm{F}_n(c_n)\)，\(x\) 是 F 的一个附着点，于是 \(x = g(p)\)，因为空间 X 是豪斯多夫的（TG, I, p. 52, 命题 1）。点 \(p\) 属于 \(\overline{\varphi_1(c_1)}\)，因而属于 \(\varphi_0(c_0)\)。于是，\(x\in \mathrm{F}_0(c_0)\)。

于是，集合 \(\mathrm{F}_{0}^{*}(c_{0}) - \mathrm{F}_{0}(c_{0})\) 包含于 G。它因而是 X 的一个贫子集，从而也是可达子集。由于 \(\mathrm{F}_{0}(c_{0}) = \mathrm{F}_{0}^{*}(c_{0}) - (\mathrm{F}_{0}^{*}(c_{0}) - \mathrm{F}_{0}(c_{0}))\) 且 \(\mathrm{F}_{0}^{*}(c_{0})\) 可达，故 \(\mathrm{F}_{0}(c_{0})\) 可达，因为 X 的可达子集构成一个 \(\sigma\)-代数。由于 \(C_{0}\) 可数，作为诸 \(\mathrm{F}_{0}(c_{0})\)（\(c_{0} \in C_{0}\)）之并的集合 S 仍是可达的。

\subsection{波兰空间中的贫集等价关系}

引理 3. --- 设 \((\mathrm{U}_{i})_{i\in\mathrm{I}}\) 是拓扑空间 X 的非空开集的一个有限族，O 是 \(X\times X\) 的一个开稠密子集。存在 X 的非空开子集的一个族 \((\mathrm{V}_{i})_{i\in\mathrm{I}}\)，使 \(V_{i}\subset U_{i}\)（对每个 \(i\in I\)），且 \(V_{i}\times V_{i'}\subset O\)（对 I 中两个不同元素的每个偶 \((i,i')\)）。

对 I 中两个不同元素的每个偶 \((j_{1}, j_{2})\)，由族 \((x_{i}) \in \mathrm{X}^{\mathrm{I}}\)（满足 \((x_{j_{1}}, x_{j_{2}}) \in \mathrm{O}\)）所成的集合是 \(X^{I}\) 的一个开稠密子集。因此，这些开集之交 \(\Omega\)——当 \((j_{1}, j_{2})\) 取遍 I 的不同元素之偶所成的有限集时——是 \(X^{I}\) 的一个开稠密集。

因此，\(\Omega \cap \prod_{i\in \mathbf{I}}\mathrm{U}_i\) 是 \(\mathrm{X}^{\mathrm{I}}\) 的一个非空开集，从而包含 \(\mathrm{X}^{\mathrm{I}}\) 的一个形如 \(\prod_{i\in \mathbf{I}}\mathrm{V}_i\) 的开集，其中对每个 \(i\in \mathbf{I}\)，\(\mathrm{V}_i\) 是 \(\mathrm{U}_i\) 的一个非空开子集。这就证明了引理。

命题 8. --- 设 P 是一个非空波兰拓扑空间，R 是 P 上的一个等价关系，其图象是 P × P 的一个贫子集。存在一个从集合 \(\{0,1\}^{\mathbf{N}}\)（赋以诸离散拓扑的乘积拓扑）到 P 中的连续单射映射，其像与 R 的每个等价类至多交于一点。

给空间 P 赋以一个与其拓扑相容且使 P 完备的度量 \(d\)。设 \((\mathrm{A}_n)_{n\in \mathbf{N}}\) 是 \(\mathrm{P}\times \mathrm{P}\) 的闭子集的一个递增序列，其内部皆为空，使得关系 R 的图象 \(\Gamma_{\mathrm{R}}\) 包含于诸 \(\mathrm{A}_n\) 之并。

设 \(\mathcal{O}\) 是 P 的非空开子集所成的集合。我们将用归纳法构造一个序列 \((f_n)_{n \in \mathbf{N}}\)，其中对每个 \(n\)，\(f_n\) 是从 \(\{0,1\}^n\) 到 \(\mathcal{O}\) 中的一个映射，它满足下列性质：

\begin{enumerate}
\def\labelenumi{(\roman{enumi})}
\item
  对每个 \(n \geqslant 1\)、每个 \(x \in \{0,1\}^{n-1}\) 及每个 \(t \in \{0,1\}\)，开集 \(f_n(x,t)\) 的闭包包含于 \(f_{n-1}(x)\)；
\item
  对每个 \(x \in \{0,1\}^n\)，有 \(\mathrm{diam}(f_n(x)) \leqslant 2^{-n}\)；
\item
  对每个 \(n \geqslant 1\) 及两个不同元素的每个偶 \((x, x')\)（属于 \(\{0, 1\}^n\)），\(f_n(x) \times f_n(x')\) 与 \(\mathrm{A}_{n-1}\) 不相交。
\end{enumerate}

选取 P 的一个直径 \(\leqslant 1\) 的非空开集 U，并把 \(f_{0}\) 定义为以 \(\{U\}\) 为像的常值映射。假设映射 \(f_{0}, \ldots, f_{n}\) 已经构造好。

设 \(p \colon \{0,1\}^{n+1} \to \{0,1\}^n\) 是由 \(p(x_0, \ldots, x_n) = (x_0, \ldots, x_{n-1})\) 定义的映射。把上述引理 3 应用于开集族 \((f_n(p(x)))\)（对 \(x \in \{0,1\}^{n+1}\)）及开稠密集 \(\mathbb{C}\mathrm{A}_n\)（\(\mathrm{P} \times \mathrm{P}\) 的开稠密集），可知存在 P 的非空开子集的一个族 \((g(x))_{x \in \{0,1\}^{n+1}}\)，使 \(g(x) \subset f_n(p(x))\)（对所有 \(x\)），且 \((g(x) \times g(x')) \cap \mathrm{A}_n\) 为空（对两个不同元素的每个偶 \((x,x')\)，属于 \(\{0,1\}^{n+1}\)）。然后定义映射 \(f_{n+1}\) 如下：对每个元素 \(x \in \{0,1\}^{n+1}\)，选取 \(g(x)\) 的一个非空开子集，其直径 \(\leqslant 2^{-n-1}\) 且其闭包包含于 \(g(x)\)。

对每个元素 \(x = (x_n)_{n \in \mathbf{N}}\)（属于 \(\{0,1\}^{\mathbf{N}}\)），集合的序列 \((f_n(x_0, \ldots, x_{n-1}))_{n \in \mathbf{N}}\) 是 P 的开子集的一个递减序列，其中每一个都包含后一个的闭包，且其直径趋于 0；因此这个集合序列的交退缩为一点（TG, II, p. 15），把它记作 \(f(x)\)。若两点 \(x\)、\(x'\)（属于 \(\{0,1\}^{\mathbf{N}}\)）满足 \(x_i = x'_i\)（对 \(i \leqslant n\)），则有 \(d(f(x), f(x')) \leqslant 2^{-n}\)。于是，映射 \(f: \{0,1\}^{\mathbf{N}} \to \mathbf{P}\) 连续。设 \(x\) 与 \(x'\) 是 \(\{0,1\}^{\mathbf{N}}\) 的两个不同元素。对 \(n \in \mathbf{N}\)（使 \((x_0, \ldots, x_n) \neq (x_0', \ldots, x_n')\)），开集 \(f_{n+1}(x_0, \ldots, x_n) \times f_{n+1}(x_0', \ldots, x_n')\) 与 \(A_n\) 不交（按 \(f_{n+1}\) 的定义），故偶 \((f(x), f(x'))\) 不属于 \(A_n\)。由此得出 \(f(x)\) 与 \(f(x')\) 对关系 R 不等价。于是，\(f\) 是单射的，且其像与 R 的每个等价类至多交于一点。命题由此得证。

\subsection{庞加莱群的基数}

命题 9. --- 设 X 是一个可解圈拓扑空间，W 是 X 的拓扑的一个基。对 X 的每一点 x，群 \(\pi_{1}(X,x)\) 的基数以 sup(Card(W),Card(N)) 为上界。特别地，可解圈的可数型度量空间的庞加莱群是可数的。

事实上，空间 \(\lambda_{x}(\mathrm{X})\)——赋以终点映射——是 X 的一个连通覆叠，它在 x 处的纤维是 \(\pi_{1}(\mathrm{X}, x)\)。断言于是由 I, p. 40, 定理 3 得出。

引理 4. --- 设 X 是贝尔空间，G 是连续地作用在 X 上的一个拓扑群，B 是 X 的一个不贫的可达子集。由点 \(g \in G\)（满足 \(B \cap gB \neq \emptyset\)）所成的集合是 G 的单位元的一个邻域。

由于 B 可达，\(\complement B\) 也可达，因为 X 的可达子集的全体是一个 \(\sigma\)-代数。于是存在 X 的一个开子集 U，使 \(U \cap \complement B\) 与 \(B \cap \complement U\) 都在 X 中贫。

设 V 是 G 的单位元的一个邻域，W 是 X 的一个包含于 U 的非空开子集，使 \(V \cdot W \subset U\)。对所有 \(g \in V\)，\(U \cap gU\) 非空。

设 \(g \in G\) 使 \(B \cap gB\) 为空。关系式

\[
\begin{array}{l} \mathrm{U} \cap g \mathrm{U} = (\mathrm{U} \cap g \mathrm{U}) \cap \mathbb {C} (\mathrm{B} \cap g \mathrm{B}) \\ \qquad = (\mathrm{U} \cap g \mathrm{U}) \cap (\mathbb {C} \mathrm{B} \cup g \mathbb {C} \mathrm{B}) \\ \qquad = (\mathrm{U} \cap g \mathrm{U} \cap \mathbb {C} \mathrm{B}) \cup (\mathrm{U} \cap g \mathrm{U} \cap g \mathbb {C} \mathrm{B}) \\ \qquad \subset (\mathrm{U} \cap \mathbb {C} \mathrm{B}) \cup g (\mathrm{U} \cap \mathbb {C} \mathrm{B}) \end{array}
\]

蕴含 U ∩ gU 在 X 中贫。由于它是一个开子集且 X 是贝尔空间，它为空。于是有 \(g \notin V\)，引理由此得证。

定理 2（Shelah \(^{(2)}\)）. --- 设 X 是一个连通且局部道路连通的波兰空间，x 是 X 的一点。若 X 不可解圈，则群 \(\pi_{1}(X,x)\) 具有连续统的基数。

设 d 是定义 X 的拓扑且使 X 完备的一个度量。假设 X 不可解圈。则存在一点 \(a \in X\)，且对每个整数 \(n \geqslant 0\)，存在 X 中的一个圈 \(c_{n}\)，其像的直径 \(\leqslant 2^{-n}\)，且其类在 \(\pi_{1}(X, a)\) 中非平凡。

设 K 表示集合 \(\{0,1\}^{\mathbf{N}}\)，并给它赋以乘积拓扑——空间 \(\{0,1\}\) 赋以离散拓扑。对 K 的每个元素 \(\varepsilon = (\varepsilon_n)\)，设 \(c_{\varepsilon}\) 是由 \(c_{\varepsilon}(0) = a\) 定义的、从 I 到 X 中的映射，且对 \(2^{-n - 1} \leqslant t \leqslant 2^{-n}\)，有 \(c_{\varepsilon}(t) = c_n(2^{n + 1}t - 1)\)（若 \(\varepsilon_n = 1\)），否则 \(c_{\varepsilon}(t) = a\)。我们有 \(c_{\varepsilon}(0) = c_{\varepsilon}(1) = a\)。映射 \(c_{\varepsilon}\) 在每一点 \(t \in ]0,1]\) 处连续。它在 0 处也连续，因为有 \(d(a,c_{\varepsilon}(t)) \leqslant 2^{-n}\)（若 \(t \in [0,2^{-n}] \)）。于是映射 \(c_{\varepsilon}\) 是 X 中以 \(a\) 为基点的一个圈。

若 \(\varepsilon\) 与 \(\varepsilon'\) 是 K 的元素，满足 \((\varepsilon_0, \ldots, \varepsilon_n) = (\varepsilon_0', \ldots, \varepsilon_n')\)，则 \(c_{\varepsilon}(t) = c_{\varepsilon'}(t)\)（对所有 \(t \in [2^{-n-1}, 1]\)），且有 \(d(c_{\varepsilon}(t), c_{\varepsilon'}(t)) \leqslant d(a, c_{\varepsilon}(t)) + d(a, c_{\varepsilon'}(t)) \leqslant 2^{-n}\)（若 \(t \in [0, 2^{-n-1}]\)）。由此得出，映射 \(\varepsilon \mapsto c_{\varepsilon}\)（从 K 到空间 \(\Omega_a(X)\) 中）是连续的，当空间 \(\Omega_a(X)\) 赋以紧收敛拓扑时。

记 \(\Gamma \subset \mathrm{K} \times \mathrm{K}\) 为由满足下列条件的偶 \((\varepsilon, \varepsilon')\) 所成的集合：\(c_{\varepsilon}\) 严格同伦于 \(c_{\varepsilon'}\)。它是 K 上一个等价关系 R 的图象。

引理 5. --- 集合 \(\Gamma\) 是 K × K 的一个贫子集。

设 Z 表示空间 \(K \times K \times \mathcal{C}_{c}(\mathbf{I} \times \mathbf{I}; X)\)。空间 \(\mathcal{C}_{c}(\mathbf{I} \times \mathbf{I}; X)\) 的拓扑由度量 \(\delta\)——由 \(\delta(h, h') = \sup_{u \in \mathbf{I} \times \mathbf{I}} d(h(u), h'(u))\) 给出——定义，且空间对这个度量是完备的（TG, X, p. 20, 推论及 TG, X, p. 9, 推论 1）；由 TG, X, p. 24, 定理 1，这个拓扑是可数型的。于是空间 \(\mathcal{C}_{c}(\mathbf{I} \times \mathbf{I}; X)\) 是波兰空间。空间 Z 亦然，因为 K 是波兰空间（TG, IX, p. 57, 命题 1）。

设 H 是由满足下列条件的三元组 \((\varepsilon,\varepsilon^{\prime},h)\) 组成的 Z 的子集：h 是连接 \(c_{\varepsilon}\) 与 \(c_{\varepsilon^{\prime}}\) 的一个严格同伦。从 Z 到 \(X^{2}\) 中的映射——由 \(a_{t}\colon(\varepsilon,\varepsilon^{\prime},h)\mapsto(c_{\varepsilon}(t),h(t,0))\)、\(b_{t}\colon(\varepsilon,\varepsilon^{\prime},h)\mapsto(c_{\varepsilon^{\prime}}(t),h(t,1))\) 与 \(c_{t}\colon(\varepsilon,\varepsilon^{\prime},h)\mapsto(h(0,t),h(1,t))\) 给出——对每个 \(t\in I\) 都连续，因为映射 \(\varepsilon\mapsto c_{\varepsilon}\)（从 K 到 \(\Omega_{a}(X)\)）连续，映射 \(h\mapsto h(s,t)\)（从 \(\mathcal{C}_{c}(\mathbf{I}\times\mathbf{I};\mathbf{X})\) 到 X 中）亦然。按定义，H 是诸集合 \(a_{t}^{-1}(\Delta_{\mathrm{X}})\)、\(b_{t}^{-1}(\Delta_{\mathrm{X}})\) 与 \(c_{t}^{-1}((a,a))\)（对 \(t\in I\)）之交，其中 \(\Delta_{X}\) 表示 X 的对角线。这就证明了 H 是 Z 的一个闭子集。

于是，H 是波兰空间。设 \(p \colon \mathrm{Z} \to \mathrm{K} \times \mathrm{K}\) 是典范投影；按定义我们有 \(\Gamma = p(\mathrm{H})\)。由于 \(\mathrm{K} \times \mathrm{K}\) 是分离的，\(\Gamma\) 是 \(\mathrm{K} \times \mathrm{K}\) 的一个苏斯林子空间（TG, IX, p. 59, 定义 2）。由 IV, p. 362, 命题 7，它是 \(\mathrm{K} \times \mathrm{K}\) 的一个可达子集。

假设 \(\Gamma\) 不贫。空间 K × K 是紧拓扑空间，因而是贝尔空间（TG, IX, p. 55, 定理 1）。给群 \(G_{0}\)（集合 \(\{0,1\}\) 的置换所成的群）赋以离散拓扑，并让乘积拓扑群 \(G = G_{0}^{\mathbf{N}}\) 对角地作用在 \(K = \{0,1\}^{\mathbf{N}}\) 上。于是群 G 通过映射 (g, (x,y)) ↦ (x, g · y) 连续地作用在 K × K 上。由引理 4，由元素 \(g \in G\)（满足 \(\Gamma \cap g \cdot \Gamma \neq \varnothing\)）所成的集合 V 是 G 的单位元的一个邻域。

设 \(g \in V\)；设 \((\varepsilon, \varepsilon') \in \Gamma \cap g \cdot \Gamma\)。由于 \((\varepsilon, \varepsilon') \in \Gamma\)，有 \(\varepsilon\,R\,\varepsilon'\)；由于 \((\varepsilon, \varepsilon') \in g \cdot \Gamma\)，有 \(g^{-1} \cdot (\varepsilon, \varepsilon') = (\varepsilon, g^{-1}\cdot\varepsilon') \in \Gamma\)，因而 \(\varepsilon\,R\,g^{-1}\cdot\varepsilon'\)。我们这样就证明了：对每个 \(g \in V\)，存在 \(\varepsilon \in K\)，使 \(\varepsilon\) 与 \(g\varepsilon\) 对 R 等价。

对 \(m \in \mathbf{N}\)，记 \(\tau_m\) 为 G 的这样的元素：它的所有项都等于 \(e\)，只有指标为 \(m\) 的那一项等于非平凡元素 \(\tau\)（属于 \(G_0\)）。存在一个整数 \(m\) 使 \(\tau_m\) 属于 V；于是选取 \(\varepsilon \in K\)，使 \(\varepsilon\) 与 \(\varepsilon' = \tau_m \cdot \varepsilon\) 对 R 等价。这蕴含圈 \(c_{\varepsilon}\) 与 \(c_{\varepsilon'}\) 严格同伦。按构造，这两个圈在区间 \([0, 2^{-m-1}]\) 与 \([2^{-m}, 1]\) 上一致；在区间 \([2^{-m-1}, 2^{-m}]\) 上，一个是像为 a 的常值映射，另一个是映射 \(t \mapsto c_{m}(2^{m+1}t - 1)\)。由此得出 \(c_{m}\) 严格同伦于像为 \(\{a\}\) 的常值圈，从而得到矛盾。引理 5 由此得证。

我们现在来完成定理 2 的证明。由命题 8，存在一个连续单射映射 \(\gamma\)（从 \(\{0,1\}^{\mathbf{N}}\) 到 K 中），其像与 R 的每个等价类至多交于一点。若 \(k\) 与 \(k'\) 是 \(\{0,1\}^{\mathbf{N}}\) 的两个不同元素，则圈 \(c_{\gamma(k)}\) 与 \(c_{\gamma(k')}\) 在 X 中不严格同伦，且映射 \(\{0,1\}^{\mathbf{N}} \to \pi_1(\mathrm{X},a)\)（由 \(k \mapsto [c_{\gamma(k)}]\) 给出）是单射的。特别地，\(\operatorname{Card}(\pi_1(\mathrm{X},a)) \geqslant \operatorname{Card}(\{0,1\}^{\mathbf{N}}) = \operatorname{Card}(\mathfrak{P}(\mathbf{N}))\)。由于 X 是可数型的可度量化拓扑空间，\(\Omega_a(\mathrm{X})\) 亦然（TG, X, p. 24, 定理 1）。于是，\(\Omega_a(\mathrm{X})\) 同胚于 \([0,1]^{\mathbf{N}}\) 的一个子空间（TG, IX, p. 18, 命题 12），且

\[
\begin{array}{r l} & {\mathrm{Card} (\Omega_ {a} (\mathbf {X})) \leqslant \mathrm{Card} ([ 0, 1 ] ^ {\mathbf {N}}) = \mathrm{Card} (\mathfrak {P} (\mathbf {N}) ^ {\mathbf {N}})} \\ & {\qquad = \mathrm{Card} (\mathfrak {P} (\mathbf {N} \times \mathbf {N})) = \mathrm{Card} (\mathfrak {P} (\mathbf {N})).} \end{array}
\]

更有，\(\operatorname{Card}(\pi_{1}(\mathrm{X},a))\leqslant\operatorname{Card}(\mathfrak{P}(\mathbf{N}))\)。于是由 E, III, p. 25 的推论 2 得出，\(\pi_{1}(\mathrm{X},a)\) 具有连续统的基数，此即所要证明的。

例. --- 设 P 是这样的拓扑空间：它是诸圆周——以 \((2/n,0)\) 为中心、过平面 \(R^{2}\) 的原点（\(n \geqslant 1\)）——之并（III, p. 336, 习题 6）。P 的庞加莱群具有连续统的基数（III, p. 338, 习题 9）。

\section{拓扑群的庞加莱群}

\subsection{群的局部同态的扩张}

定义 1. --- 给定一个拓扑群 G 与一个群 G'，G 到 G' 中的局部同态是从 G 的单位元的一个邻域 V 到 G' 中的一个映射 f，使得对 V 中满足 xy ∈ V 的每对点 x、y，有 f(xy) = f(x)f(y)。

若 G' 是拓扑群且 f 是连续映射，则称 f 是连续局部同态（或拓扑群的局部态射）。

若 G 与 G' 是拓扑群，从 G 到 G' 的局部同构的概念已在 TG, III, p. 6, 定义 2 中定义。G 到 G' 的一个局部同构是从 G 的单位元的一个邻域 V 到 G' 的单位元的一个邻域 V' 上的一个同胚 f，使得 f 与 f 的逆映射都是局部同态。

命题 1. --- 设 G 与 G' 是拓扑群，p: G → G' 是群同态。为使 p 使 G 成为 G' 的一个覆叠，必要且充分的是：p 在 G 的单位元的某个适当邻域上的限制是 G 到 G' 的一个局部同构。

该条件由 TG, III, p. 6 的命题 3 是必要的。反之，假设 \(p\) 诱导一个从 G 的单位元 \(e\) 的开邻域 V 到 G' 的单位元的一个邻域 V' 上的同胚。于是映射 \(p\) 连续（TG, III, p. 15, 命题 23）且开（TG, III, p. 16, 命题 24）。\(p\) 的像 H 是 G′ 的一个包含 V 的子群，故在 G′ 中既开又闭（TG, III, p. 7, 推论）。设 N 是 \(p\) 的核；我们有 N∩V = \{e\}；因此 N 离散（同前， 命题 5）。于是，\(p\) 使 G 成为 H 的一个覆叠，以群 N 为主（I, p. 100, 定理 1 的推论 3）。由于 H 在 G' 中既开又闭，G'-空间 (G, p) 是一个覆叠。

命题 2. --- 设 G 是连通拓扑群，G' 是群，\(f: V \to G'\) 是 G 到 G' 中的局部同态，其中 V 是 G 的单位元的一个连通邻域。则存在一个连通拓扑群 H、一个拓扑群的态射 \(p\colon\mathrm{H}\to\mathrm{G}\)（使 \((\mathrm{H},p)\) 是 G 的一个覆叠）以及一个群同态 \(\varphi\colon\mathrm{H}\to\mathrm{G}'\)，使得由 \(y\in p^{-1}(\mathrm{V})\)（满足 \(f(p(y))=\varphi(y)\)）所成的集合是 H 中单位元的一个邻域。

若 G' 是拓扑群且映射 f 连续，则这样的同态 \(\varphi\) 是连续的。

引理 1. — 设 G 是拓扑群，V 是 G 的单位元 e 的一个连通邻域。对 e 在 G 中的每个邻域 U 及每个 \(x \in V\)，存在一个整数 \(n \in N\) 及 U 中的元素 \(u_{1}, \ldots, u_{n}\)，使 \(u_{1} \ldots u_{n} = x\)，且 \(u_{1} \ldots u_{k} \in V\)（对每个满足 \(1 \leqslant k \leqslant n\) 的整数 k）。

我们可以假设 U 开且包含于 V。设 A 表示满足引理条件的 \(x \in V\) 所成的集合。若 \(x \in A\) 且 \(y \in xU \cap V\)，则 \(y \in A\)，由此得 \(AU \cap V \subset A\)；这表明 A 在 V 中开。设 \(x \in V\) 使 \(xU^{-1} \cap A \neq \varnothing\)；于是有 \(x \in AU \cap V\)，因此 \(x \in A\)。于是，若 \(x \in V\) 且 \(x \notin A\)，则 \(xU^{-1} \cap A = \varnothing\)，故 A 在 V 中闭。由于 \(e \in A\) 且 V 连通，得出 A = V。

现在来证明命题。记 \(j\) 为映射 \(g \mapsto (g, f(g))\)（从 V 到 \(\mathbf{G} \times \mathbf{G}'\)），并设 H 是 \(\mathbf{G} \times \mathbf{G}'\) 的一个子群（由 \(j(\mathrm{V})\) 生成）。

设 U 是 \(e\) 在 G 中的一个邻域，包含于 V。设 \(x \in V\)；由引理 1，存在 \(u_1, \ldots, u_n \in U\)，使 \(x = u_1 \ldots u_n\)，且 \(u_1 \ldots u_k \in V\)（对每个整数 \(k\)，满足 \(1 \leqslant k \leqslant n\)）。由归纳法，有 \(f(u_1 \ldots u_k) = f(u_1) \ldots f(u_k)\)（对每个整数 \(k\)，\(1 \leqslant k \leqslant n\)）。特别地，\(j(x)\) 属于由 \(j(U)\) 生成的子群。由此得出 H 由 \(j(U)\) 生成。

设 \(\mathcal{B}\) 是 H 的形如 \(j(\mathrm{U})\) 的子集所成的集合，其中 U 是 \(e\) 在 V 中的一个邻域。我们来证明 H 上存在唯一的拓扑，它与 H 的群结构相容，且使 \(\mathcal{B}\) 是单位元的邻域滤子的基。为此，只需证明集合 \(\mathcal{B}\) 满足 III, p. 4 中的条件 \((\mathrm{GV}_{\mathrm{I}}^{\prime})\)、\((\mathrm{GV}_{\mathrm{II}}^{\prime})\) 与 \((\mathrm{GV}_{\mathrm{III}}^{\prime})\)。

设 U 是 e 在 G 中的一个邻域，包含于 V；存在 e 的一个邻域 \(U'\)，使 \(U' \cdot U' \subset U\)。对 \(U'\) 的每对点 x、y，有 \(xy \in V\) 且 \(f(xy) = f(x)f(y)\)。于是，\(j(\mathrm{U}') \in \mathcal{B}\) 且 \(j(\mathrm{U}') \cdot j(\mathrm{U}') \subset j(\mathrm{U})\)。这表明条件 \((\mathrm{GV}_{\mathrm{I}}')\) 满足。

于是集合 \(U'' = V \cap U^{-1}\) 是 e 在 V 中的一个邻域，且对 \(x \in U''\)，有 \(x^{-1} \in U\) 及 \(f(x^{-1}) = f(x)^{-1}\)。因而 \(j(\mathrm{U}'')^{-1} \subset j(\mathrm{U})\)，这表明条件 \((\mathrm{GV}_{\mathrm{II}}' )\) 成立。

最后，取定 e 在 V 中的一个邻域 U，使 \(U = U^{-1}\) 且 \(U^{3} \subset V\)。设 W 是 e 在 U 中的一个邻域，设 \(h = (g, f(g))\) 是 \(j(\mathrm{U})\) 的一个元素。存在 e 的一个包含于 W 的邻域 \(W'\)，使 \(gW'g^{-1} \subset W\)。于是，\(j(\mathrm{W}') \in \mathcal{B}\) 且 \(hj(\mathrm{W}')h^{-1} \subset j(\mathrm{W})\)，因为有 \(f(gxg^{-1}) = f(g)f(x)f(g^{-1})\)（对 \(x \in W'\)）。

设 \(h \in H\)。由于 U 是 e 的包含于 V 的邻域，\(j(\mathrm{U})\) 生成 H；由于 U 对称，存在 U 中的元素 \(u_{1}, \ldots, u_{n}\)，使 \(h = j(u_{1}) \ldots j(u_{n})\)。对 n 用归纳法，存在 e 的一个包含于 W 的邻域 \(W'\)，使 \(hj(\mathrm{W}')h^{-1} \subset j(\mathrm{W})\)。于是条件 \((\mathrm{GV}_{\mathrm{III}}')\) 满足。

然后给群 H 赋以这个拓扑。

记 \(p \colon \mathrm{H} \to \mathrm{G}\) 为第一个投影 \(\mathrm{G} \times \mathrm{G}' \to \mathrm{G}\) 在 H 上的限制。这是一个群同态。对 \(e\) 的每个包含于 V 的邻域 U，集合 \(p^{-1}(\mathrm{U})\) 包含 H 的单位元的邻域 \(j(\mathrm{U})\)，故 \(p\) 是拓扑群的一个连续同态。对 \(e\) 的每个包含于 V 的邻域 U，有 \(p(j(\mathrm{U})) = \mathrm{U}\)，因此 \(p\) 是开映射。由于 G 连通，\(p\) 是满射的。它的核是离散的，因为它除单位元外与 \(j(\mathrm{V})\) 不相交。于是由推论 3（I, p. 100）得出，G-空间 \((\mathrm{H}, p)\) 是一个覆叠。

设 \(\varphi\) 是第二个投影 \(G \times G' \to G'\) 在 H 上的限制。若 \(g \in V\)，有 \((g, f(g)) = j(g) \in j(V)\) 且 \(\varphi(g, f(g)) = f(g)\)，因而有 \(\varphi(y) = f(p(y))\)（对 \(y \in j(V)\)）。

进而假设 G' 是拓扑群且映射 f 在 e 处连续。于是同态 \(\varphi\) 在 H 的单位元处连续，因而连续（TG, III, p. 15, 命题 23）。

推论 1. --- 设 G 是单连通拓扑群，G' 是群。设 V 是 G 中单位元的一个连通邻域，f: V → G' 是一个局部同态。存在唯一的群同态 h: G → G' 延拓 f。若 G' 是拓扑群且映射 f 连续，则同态 h 连续。

群 G 连通。设 H、\(\varphi\) 与 p 如命题 2 中所示。由于 G 单连通，H 是 G 的一个可平凡化覆叠（I, p. 124, 定义 3）；由于 H 连通且非空，映射 \(p\) 是拓扑群的一个同构。置 \(h = \varphi \circ p^{-1}\)；映射 \(h\) 是群同态，且存在 G 的单位元 \(e\) 的一个包含于 V 的开邻域 U，使 \(f|U = h|U\)。由引理 1 得出，\(f\) 与 \(h\) 在 V 上一致。换句话说，映射 \(h\) 延拓映射 \(f\)。

若 G' 是拓扑群且映射 f 连续，则同态 \(\varphi\) 连续，故 h 也连续。

我们来证明这种延拓的唯一性。G 中使从 G 到 G′ 的两个同态一致的点所成的集合是 G 的一个子群。由于群 G 连通，G 的每个包含单位元的一个邻域的子群都等于 G（TG, III, p. 8, 命题 6）。h 的唯一性由此得出。

注 1. --- 当 G = R 时，这个推论由 TG, V, p. 3 的命题 6 得出。

推论 2. --- 两个局部连通且单连通的拓扑群若是局部同构的，则是同构的。

设 G 与 G' 是局部连通拓扑群，V 与 V' 分别是 G 与 G' 的单位元的连通邻域，\(f: V \to V'\) 是一个同胚，且是 G 到 G' 的一个局部同构。由推论 1，存在唯一的连续群同态 \(\varphi: G \to G'\) 延拓 f，且存在唯一的连续群同态 \(\varphi': G' \to G\) 延拓 \(f^{-1}\)。同态 \(\varphi' \circ \varphi\) 与 \(\operatorname{Id}_{\mathrm{G}}\) 在 e 于 G 中的一个邻域上一致，故相等，因为 G 连通。同样，\(\varphi \circ \varphi' = \operatorname{Id}_{\mathrm{G}'}\)，推理由此得证。

推论 3. --- 设 G 是单连通拓扑群，V 是 G 的单位元的一个连通邻域。以集合 V 为生成集，并以关系元集 \(\mathbf{r}\) 为诸 \(xyz^{-1}\)（其中 \((x,y,z)\) 取遍 V 中满足 \(xy = z\) 的元素三元组）所成的族，便定义了 G 的一个表现。

设 \(\mathrm{F}(\mathrm{V}, \mathbf{r})\) 是自由群 \(\mathrm{F}(\mathrm{V})\) 关于包含诸元素 \(xyz^{-1}\)（其中 \((x, y, z) \in \mathrm{V} \times \mathrm{V} \times \mathrm{V}\) 且 xy = z）的最小正规子群的商群（A, I, p. 86）。记 \(f: V \to F(V, r)\) 与 \(g: F(V, r) \to G\) 为典范映射。按构造，映射 f 是 G 到 \(\mathrm{F}(V, \mathbf{r})\) 的一个局部同态。由推论 1，存在唯一的群同态 \(\overline{f}: G \to F(V, r)\) 延拓 f。

由于群 \(F(V, \mathbf{r})\) 由 \(f(V)\) 生成，同态 \(\overline{f}\) 是满射的。对 \(x \in V\)，有 \(g(\overline{f}(x)) = g(f(x)) = x\)；由于 V 生成 G，\(g \circ \overline{f} = \operatorname{Id}_{\operatorname{G}}\)，这证明 \(\overline{f}\) 是单射的。因此它是一个同构。
